% Artigo em português; a edição ABNT carrega esta mesma fonte.
\ifdefined\besovabntedition
  \documentclass[12pt,a4paper]{article}
\else
  \documentclass[11pt,a4paper]{article}
\fi
\usepackage{conjecturas_lp_estilo}
\abntentry{artigo_1-devore1998}{DeVore}{DeVore}{1998}
\abntentry{artigo_1-devore2021}{DeVore; Hanin; Petrova}{DeVore, Hanin\articleand{}Petrova}{2021}
\abntentry{artigo_2-devorelorentz1993}{DeVore; Lorentz}{DeVore\articleand{}Lorentz}{1993}
\abntentry{artigo_2-devorepopov1988}{DeVore; Popov}{DeVore\articleand{}Popov}{1988}
\abntentry{artigo_2-hristovivanov1990}{Hristov; Ivanov}{Hristov\articleand{}Ivanov}{1990}
\abntentry{artigo_2-petrushev1988}{Petrushev}{Petrushev}{1988}
\abntentry{artigo_2-whitney1957}{Whitney}{Whitney}{1957}
\abntentry{lp-siegel2023}{Siegel}{Siegel}{2023}
\abntentry{lp-ditzian1987}{Ditzian; Totik}{Ditzian\articleand{}Totik}{1987}
\abntentry{lp-oliveiravol}{Oliveira}{Oliveira}{2026}
\abntentry{lp-ali2020}{Ali; Nouy}{Ali\articleand{}Nouy}{2020}
\abntentry{lp-yang2024}{Yang}{Yang}{2024}
\abntentry{lp-kolomoitsev2019}{Kolomoitsev; Tikhonov}{Kolomoitsev\articleand{}Tikhonov}{2019}
\abntentry{lp-ditzian2007}{Ditzian}{Ditzian}{2007}
\abntentry{lp-lean4}{Moura; Ullrich}{de Moura\articleand{}Ullrich}{2021}
\abntentry{lp-mathlib}{Mathlib Community}{The mathlib Community}{2020}


\hypersetup{
  pdftitle={Caracterização dos espaços de Besov em Lp por realização com altura aritmética controlada},
  pdfsubject={Teoremas de Jackson, espaços de Besov e estabilidade por inclusão de redes ReLU},
  pdfkeywords={espaços de Besov, aproximação em Lp, altura aritmética, funcional de realização, redes ReLU}
}

\begin{document}
\articletitle{Caracterização dos espaços de Besov em \texorpdfstring{$L^p$}{Lp} por~realização com altura aritmética controlada}
  {Pesquisador independente\ifbesovabnt
     \ifdefempty{\besovcontact}{}{\enspace\textperiodcentered\enspace\besovcontact}\fi}
  {19 de novembro de 2026}
  {A Realization Characterization of Besov Spaces in \texorpdfstring{$L^p$}{Lp} with Controlled Arithmetic Height}

\ifbesovabnt
\articleabstract{portuguese}{Resumo}{%
Caracterizamos os espaços de Besov $B^s_{p,q}([0,1])$, com $1\le p<\infty$, $0<s<2$ e $0<q\le\infty$, pela somabilidade diádica de um funcional que combina erro em $L^p$ e suavidade de segunda ordem. Os aproximantes são funções contínuas e afins por partes, com coeficientes racionais e limites simultâneos para o número de folhas da partição diádica e a altura aritmética. A estimativa direta combina projeções afins locais, uma desigualdade de Whitney e quantização nodal. Para $s<1/p$, basta uma altura da ordem de $N^{1+1/p-s}$. Para $p>1$, $s=1/p$ e $0<q<\infty$, uma agenda de altura da ordem de $(B+1)N$ dá somabilidade uniforme na bola de raio $B\ge1$ e caracteriza a pertinência ao espaço de realização nessa bola. A caracterização global sob as agendas gerais se transfere a famílias aninhadas de redes ReLU que incluem a biblioteca racional após reescalonamento linear do orçamento ou a aproximam com tolerância somável na escala de Besov. Para a taxa uniforme $N^{-s}$ por escala, os pisos de erro determinam a potência ótima da altura. Agendas de log iterado melhoram as condições suficientes de somabilidade nos regimes estudados. Em cada bola de raio fixo, o código explícito da biblioteca atinge a taxa de erro em $L^p$ com $O(N\log N)$ bits, ante uma cota inferior de entropia de ordem $N$.}
{Palavras-chave}{espaços de Besov; funcional de realização; altura aritmética; teorema de Jackson; redes ReLU.}

\articleabstract{english}{Abstract}{%
We characterize the Besov spaces $B^s_{p,q}([0,1])$, for $1\le p<\infty$, $0<s<2$, and $0<q\le\infty$, by the dyadic summability of a functional combining $L^p$ approximation error and second-order smoothness. The approximants are continuous piecewise affine functions with rational coefficients and simultaneous bounds on the number of leaves of the dyadic partition and on arithmetic height. The direct estimate combines local affine projections, a Whitney inequality, and nodal quantization. For $s<1/p$, height of order $N^{1+1/p-s}$ suffices. For $p>1$, $s=1/p$, and $0<q<\infty$, a height schedule of order $(B+1)N$ gives uniform summability on the ball of radius $B\ge1$ and characterizes membership in the realization space on that ball. Under the general schedules, the global characterization extends to nested ReLU network families that contain the rational library after a linear budget rescaling or approximate it with a tolerance that is summable on the Besov scale. For the uniform per-scale rate $N^{-s}$, lower bounds on the error determine the optimal power of the height. Iterated-logarithm schedules improve sufficient summability conditions in the regimes studied. On each ball of fixed radius, the explicit library code attains the $L^p$ error rate with $O(N\log N)$ bits, compared with an entropy lower bound of order $N$.}
{Keywords}{Besov spaces; realization functional; arithmetic height; Jackson theorem; ReLU networks.}

\medskip\noindent{\small
Versão de 19 de novembro de 2026.\par\noindent
\ifdefempty{\besovsubmitted}{}{Submissão: \besovsubmitted.\par}%
\ifdefempty{\besovaccepted}{}{Aprovação: \besovaccepted.\par}%
\par}
\else
\articleabstract{portuguese}{Resumo}{%
Caracterizamos os espaços de Besov $B^s_{p,q}([0,1])$, com $1\le p<\infty$, $0<s<2$ e $0<q\le\infty$, pela somabilidade diádica de um funcional que combina erro em $L^p$ e suavidade de segunda ordem. Os aproximantes são funções contínuas e afins por partes, com coeficientes racionais e limites simultâneos para o número de folhas da partição diádica e a altura aritmética. A estimativa direta combina projeções afins locais, uma desigualdade de Whitney e quantização nodal. Para $s<1/p$, basta uma altura da ordem de $N^{1+1/p-s}$. Para $p>1$, $s=1/p$ e $0<q<\infty$, uma agenda de altura da ordem de $(B+1)N$ dá somabilidade uniforme na bola de raio $B\ge1$ e caracteriza a pertinência ao espaço de realização nessa bola. A caracterização global sob as agendas gerais se transfere a famílias aninhadas de redes ReLU que incluem a biblioteca racional após reescalonamento linear do orçamento ou a aproximam com tolerância somável na escala de Besov. Para a taxa uniforme $N^{-s}$ por escala, os pisos de erro determinam a potência ótima da altura. Agendas de log iterado melhoram as condições suficientes de somabilidade nos regimes estudados. Em cada bola de raio fixo, o código explícito da biblioteca atinge a taxa de erro em $L^p$ com $O(N\log N)$ bits, ante uma cota inferior de entropia de ordem $N$.}
{Palavras-chave}{Espaços de Besov, funcional de realização, altura aritmética, teorema de Jackson, redes ReLU.}

\articleabstract{english}{Abstract}{%
We characterize the Besov spaces $B^s_{p,q}([0,1])$, for $1\le p<\infty$, $0<s<2$, and $0<q\le\infty$, by the dyadic summability of a functional combining $L^p$ approximation error and second-order smoothness. The approximants are continuous piecewise affine functions with rational coefficients and simultaneous bounds on the number of leaves of the dyadic partition and on arithmetic height. The direct estimate combines local affine projections, a Whitney inequality, and nodal quantization. For $s<1/p$, height of order $N^{1+1/p-s}$ suffices. For $p>1$, $s=1/p$, and $0<q<\infty$, a height schedule of order $(B+1)N$ gives uniform summability on the ball of radius $B\ge1$ and characterizes membership in the realization space on that ball. Under the general schedules, the global characterization extends to nested ReLU network families that contain the rational library after a linear budget rescaling or approximate it with a tolerance that is summable on the Besov scale. For the uniform per-scale rate $N^{-s}$, lower bounds on the error determine the optimal power of the height. Iterated-logarithm schedules improve sufficient summability conditions in the regimes studied. On each ball of fixed radius, the explicit library code attains the $L^p$ error rate with $O(N\log N)$ bits, compared with an entropy lower bound of order $N$.}
{Keywords}{Besov spaces, realization functional, arithmetic height, Jackson theorem, ReLU networks.}
\fi

\section{Introdução}
\label{lp:sec:intro}

A teoria da aproximação caracteriza a regularidade de Besov por meio de erros associados a famílias de aproximantes, entre elas os espaços de splines e as classes de funções racionais \citep{artigo_1-devore1998,artigo_2-devorepopov1988,artigo_2-devorepopov1988b,artigo_2-petrushev1988,lp-petrushevpopov1987}. A caracterização aqui apresentada se distingue desses antecedentes pela restrição simultânea sobre a altura aritmética dos coeficientes, incorporada ao funcional de realização e quantificada pelas agendas de altura. Os funcionais de realização incorporam ao erro de aproximação uma medida de suavidade do aproximante \citep{lp-ditzian2007,artigo_2-hristovivanov1990,lp-kolomoitsev2019}. As estimativas para aproximação de classes de Besov por redes de unidades lineares retificadas (ReLU) vinculam essa regularidade ao orçamento das redes \citep{lp-ali2020,artigo_1-devore2021,lp-siegel2023,lp-yang2024}.

O orçamento de aproximação controla o número de folhas da árvore diádica e a altura dos racionais que definem os segmentos afins. Para cada orçamento $N$ e agenda $H$, o funcional de realização minimiza a soma do erro em $L^p$ e do módulo de segunda ordem do aproximante na escala $N^{-1}$. Sob as agendas e os regimes de parâmetros do \cref{lp:conj:besov}, a somabilidade diádica do funcional caracteriza $B^s_{p,q}([0,1])$. A inclusão da biblioteca racional em uma família aninhada de redes, com ampliação linear de $N$, preserva o espaço caracterizado; o mesmo ocorre sob inclusão aproximada em $L^p$ com tolerância $\varepsilon_N$ somável na escala de Besov. Esses resultados estendem à norma $L^p$ a caracterização uniforme e o teorema de estabilidade da Parte~IV de \citet{lp-oliveiravol}. A \cref{lp:sec:uniform} acrescenta cotas uniformes sobre bolas, o piso aritmético e a complexidade em bits dos aproximantes, e agendas de log iterado.

O módulo de segunda ordem controla o erro do quase-interpolante construído com projeções afins locais. A quantização de seus valores nodais produz o aproximante racional. Acima de $s=1/p$, a cota de Jackson para um alvo fixo tem termo de quantização $O(H^{-1})$ sob as hipóteses do \cref{lp:hyp:jackson}. Para $0<s<1/p$, a cota em sup dos quase-interpolantes leva à altura $N^{1+1/p-s}$. Na diagonal $s=1/p$, o \cref{lp:conj:besov} dá uma agenda global da forma $N[\log(eN)]^{\theta_q+\kappa}$, com as condições de seu item (iii). Para $p>1$ e $0<q<\infty$, o controle dos saltos nodais permite usar $H_B(N)\asymp(B+1)N$ na bola de raio $B\ge1$ pelo \cref{lp:thm:critical-linear}. A desigualdade inversa vale para qualquer família de aproximantes em $L^p$.

A \cref{lp:cor:joint-necessity} fixa uma escala $N$ e uma constante de erro $C$ comum a toda a bola unitária. Para a taxa $CN^{-s}$, a altura deve ser pelo menos da ordem de $N^{1+1/p-s}$ se $s<1/p$, de $N$ se $1/p\le s\le1$ e de $N^s$ se $1<s<2$. As cotas superiores por escala atingem essas potências, com fatores logarítmicos em alguns casos da diagonal. A somabilidade entre escalas tem condições suficientes próprias, reunidas na \cref{lp:tab:schedules}.

\section{Espaços de Besov e biblioteca racional}
\label{lp:sec:defs}

\subsection{Módulos de suavidade}

Fixe $1\le p<\infty$. Escrevemos $L^p=L^p([0,1])$ e $\|\cdot\|_p=\|\cdot\|_{L^p([0,1])}$. Para $f\in L^p$ e $0<h\le1/2$, definimos
\begin{equation}
 \Delta_h^2f(x)=f(x+2h)-2f(x+h)+f(x),\qquad x\in[0,1-2h].
 \label{lp:eq:second-difference}
\end{equation}
O módulo de segunda ordem é
\begin{equation}
 \wt{f}{t}{p}
 =\sup_{0<h\le\min\{t,1/2\}}\|\Delta_h^2f\|_{L^p([0,1-2h])},
 \qquad t>0.
 \label{lp:eq:modulus-p}
\end{equation}
Cada termo da diferença é avaliado em $[0,1]$. Para $f$ contínua, escrevemos
$\omega_1(f,t)_\infty=\sup\{|f(x)-f(y)|:x,y\in[0,1],\ |x-y|\le t\}$.

Para $0<s<2$ e $0<q\le\infty$, adotamos a caracterização diádica de Besov
\citep{artigo_2-devorelorentz1993,lp-ditzian1987}
\begin{equation}
 \|f\|^*_{\Bes{s}{p}{q}}
 =\|f\|_p+\big\|\big(2^{rs}\wt{f}{2^{-r}}{p}\big)_{r\ge0}\big\|_{\ell^q},
 \qquad
 \Bes{s}{p}{q}=\{f\in L^p:\|f\|^*_{\Bes{s}{p}{q}}<\infty\}.
 \label{lp:eq:besov-p}
\end{equation}
Para $q=\infty$, usa-se o supremo. Para $s>1/p$, o \cref{lp:lem:embedding} fornece o representante contínuo de cada classe de Besov, usado nas avaliações pontuais. Para $s\le1/p$, a construção usa médias locais e dispensa avaliações pontuais de $f$.

\subsection{Altura aritmética e funcional de realização}

Para uma fração reduzida $u=a/b\in\Q$, com $b>0$, ponha
$\height(u)=\max\{|a|,b\}$. Seja $J_{k,m}=[m2^{-k},(m+1)2^{-k}]$ uma célula diádica, com centro $x_{k,m}=(m+\tfrac12)2^{-k}$. Sobre essa célula, um segmento da biblioteca $\Library$ tem a forma $u+v(x-x_{k,m})$, com $u\in\Q$ e $v\in\Q\setminus\{0\}$. Sua altura é
\[
 \max\{2,k+1,|m|+1,\height(u),\height(v)\}.
\]

Uma representação sobre $[0,1]$ é obtida por bisseções diádicas sucessivas, com um segmento em cada folha e coincidência dos valores nos extremos compartilhados. Denotamos por $\Aclass_{N,H}([0,1])$ o conjunto das funções contínuas que admitem tal representação com até $N$ folhas de altura no máximo $H$ \citep{lp-oliveiravol}.

Para uma agenda não decrescente $H\colon\N\to\N$, definimos
\begin{equation}
 \Sig{N}=\{0\}\cup\Aclass_{N,H(N)}([0,1]).
 \label{lp:eq:sigma}
\end{equation}
Com $\nu_s=\max\{1,s\}$, as agendas utilizadas são
\begin{align}
 H_{s,\kappa}(N)
 &=\left\lceil8(N+1)^{\nu_s}[\log(eN)]^\kappa\right\rceil,
 &&\kappa\ge0,\quad N\ge1,\label{lp:eq:schedule}\\
 \widetilde H_{s,\eta}(N)
 &=\left\lceil8(N+1)^{\max\{1,s+\eta\}}\right\rceil,
 &&\eta>0,\quad N\ge1.\label{lp:eq:slack}
\end{align}
Com $\theta_q=\max\{0,1-1/q\}$, de modo que $\theta_\infty=1$, as agendas dos regimes $s\le1/p$ são
\begin{align}
 \widehat H_{s}(N)
 &=\left\lceil8(N+1)^{1+1/p-s}\right\rceil,
 &&0<s<1/p,\quad N\ge1,\label{lp:eq:schedule-low}\\
 \widehat H_{q,\kappa}(N)
 &=\left\lceil8(N+1)[\log(eN)]^{\theta_q+\kappa}\right\rceil,
 &&s=1/p,\quad\kappa\ge0,\quad N\ge1.\label{lp:eq:schedule-diag}
\end{align}
A dependência de $p$ em $\widehat H_s$ é implícita.

\begin{definicao}[Funcional e espaço de realização]
\label{lp:def:realization}
Para $\mathcal G=(\mathcal G_N)_{N\ge2}$, com $\mathcal G_N\subseteq L^p$, e $f\in L^p$, o funcional de realização é
\begin{equation}
 \Rfun{\mathcal G}{N}(f)_p
 =\inf_{g\in\mathcal G_N}\bigl\{\|f-g\|_p+\wt{g}{N^{-1}}{p}\bigr\},
 \label{lp:eq:realization-p}
\end{equation}
com a convenção $\inf\varnothing=+\infty$. O espaço de realização é
\begin{equation}
 \Rstr{s}{q}{\mathcal G}_p
 =\Bigl\{f\in L^p:\big\|\big(2^{rs}\Rfun{\mathcal G}{2^r}(f)_p\big)_{r\ge1}\big\|_{\ell^q}<\infty\Bigr\}.
 \label{lp:eq:stratum-p}
\end{equation}
Para a família $(\Sig{N})_{N\ge2}$, escrevemos $\Rfun{\Library,H}{N}$ e $\Rstr{s}{q}{\Library,H}_p$.
\end{definicao}

\begin{definicao}[Família aninhada que contém a biblioteca]
\label{lp:def:network-family}
Uma família $\mathcal N=(\mathcal N_N)_{N\ge1}$ de subconjuntos de $C([0,1])$ é aninhada e contém a função nula quando
\begin{equation}
 0\in\mathcal N_N,\qquad \mathcal N_N\subseteq\mathcal N_{N+1}\qquad(N\ge1).
 \label{lp:eq:nesting}
\end{equation}
Ela contém a biblioteca sob a agenda $H$ com constante $c\ge1$ quando
\begin{equation}
 \Sig{N}\subseteq\mathcal N_{\lceil cN\rceil}\qquad(N\ge1).
 \label{lp:eq:contains}
\end{equation}
Mais geralmente, para uma sequência $(\varepsilon_N)_{N\ge1}$ em $[0,\infty)$, dizemos que $\mathcal N$ aproxima a biblioteca sob a agenda $H$ com constante $c\ge1$ e erro $(\varepsilon_N)$ em $L^p$ quando
\begin{equation}
 \forall N\ge1,\ \forall g\in\Sig{N},\ \exists\tilde g\in\mathcal N_{\lceil cN\rceil}\colon\ \|g-\tilde g\|_p\le\varepsilon_N.
 \label{lp:eq:approx-contains}
\end{equation}
\end{definicao}

\section{Caracterização de Besov e estabilidade}
\label{lp:sec:statements}

\begin{teorema}[Caracterização por realização aritmética]
\label{lp:conj:besov}
Sejam $1\le p<\infty$ e $0<s<2$.
\begin{enumerate}[label=(\roman*)]
\item Se $s>1/p$, valem as seguintes igualdades em $L^p([0,1])$:
\begin{align}
 \Bes{s}{p}{\infty}
 &=\Rstr{s}{\infty}{\Library,H_{s,0}}_p,
 &&\label{lp:eq:conj-endpoint}\\
 \Bes{s}{p}{q}
 &=\Rstr{s}{q}{\Library,H_{s,0}}_p,
 &&0<q<\infty,\quad s<1,\label{lp:eq:conj-subcritical}\\
 \Bes{s}{p}{q}
 &=\Rstr{s}{q}{\Library,H_{s,\kappa}}_p,
 &&0<q<\infty,\quad s\ge1,\quad\kappa>1/q,\label{lp:eq:conj-critical}\\
 \Bes{s}{p}{q}
 &=\Rstr{s}{q}{\Library,\widetilde H_{s,\eta}}_p,
 &&0<q<\infty,\quad\eta>0.
 \label{lp:eq:besov-tilde}
\end{align}
\item Se $0<s<1/p$ e $0<q\le\infty$, então
\begin{equation}
 \Bes{s}{p}{q}=\Rstr{s}{q}{\Library,\widehat H_{s}}_p.
 \label{lp:eq:conj-low}
\end{equation}
\item Se $s=1/p$, $0<q\le\infty$ e $\kappa\ge0$, com $\kappa>1/q$ quando $p=1$ e $q<\infty$, então
\begin{equation}
 \Bes{s}{p}{q}=\Rstr{s}{q}{\Library,\widehat H_{q,\kappa}}_p.
 \label{lp:eq:conj-diag}
\end{equation}
\end{enumerate}
\end{teorema}

\begin{teorema}[Estabilidade por inclusão de redes]
\label{lp:conj:stability}
Sejam $p,s,q,H$ como em qualquer dos casos do \cref{lp:conj:besov}. Se $\mathcal N$ satisfaz \eqref{lp:eq:nesting} e \eqref{lp:eq:contains} para algum $c\ge1$, então
\begin{equation}
 \Bes{s}{p}{q}=\Rstr{s}{q}{\mathcal N}_p.
 \label{lp:eq:network-realization-p}
\end{equation}
\end{teorema}

\begin{teorema}[Estabilidade por inclusão aproximada]
\label{lp:thm:approx-stability}
Sejam $p,s,q,H$ como em qualquer dos casos do \cref{lp:conj:besov}. Se $\mathcal N$ satisfaz \eqref{lp:eq:nesting} e \eqref{lp:eq:approx-contains} para algum $c\ge1$ e uma sequência $(\varepsilon_N)_{N\ge1}$ em $[0,\infty)$ tal que
\begin{equation}
 \big(2^{rs}\varepsilon_{2^r}\big)_{r\ge1}\in\ell^q,
 \label{lp:eq:approx-summable}
\end{equation}
então $\Bes{s}{p}{q}=\Rstr{s}{q}{\mathcal N}_p$.
\end{teorema}

Para uma família de redes ReLU, a inclusão \eqref{lp:eq:contains} é uma condição sobre as funções realizadas. Por exemplo, uma função afim por partes com nós internos $t_1<\cdots<t_{m-1}$ e inclinações consecutivas $\beta_1,\ldots,\beta_m$ admite em $[0,1]$ a representação
\[
 g(x)=g(0)+\beta_1 x+\sum_{j=1}^{m-1}(\beta_{j+1}-\beta_j)(x-t_j)_+,
 \qquad u_+=\max\{u,0\}.
\]
Como $x=x_+$ em $[0,1]$, a expressão é realizável por uma camada oculta de largura no máximo $m$. Assim, as classes de funções realizadas por redes de uma camada oculta e largura no máximo $N$ contêm $\Sig{N}$ com $c=1$. O \cref{lp:thm:approx-stability} estende essa estabilidade a arquiteturas que reproduzem a biblioteca apenas com precisão controlada em $L^p$, desde que a tolerância satisfaça \eqref{lp:eq:approx-summable}. Verificar essa tolerância para uma arquitetura concreta, como redes com pesos quantizados, é uma questão à parte, que o texto não trata.

Nas agendas dos itens (ii) e (iii), a cota dos valores nodais envolve $n^{1/p-s}$ abaixo da diagonal e $(\log n)^{\theta_q}$ na diagonal. O expoente $\theta_q$ registra esse controle dos valores. O expoente adicional $\kappa$ é usado na diagonal $p=s=1$ para somar a cota do erro de quantização quando $q<\infty$. A \cref{lp:sec:ball} trata das agendas uniformes em bolas de raio fixo.

\section{Jackson racional com altura controlada}
\label{lp:sec:program}

\begin{teorema}[Estimativa de Jackson em $L^p$]
\label{lp:hyp:jackson}
Para $1\le p<\infty$ existe $C_W=C_W(p)\ge1$ com a seguinte propriedade. Se $1/p<s<2$ e $f\in\Bes{s}{p}{\infty}$, seja $R=\|f\|_\infty$ a norma do representante contínuo. Existe $n_0=n_0(f,p,s)\ge2$ tal que, para toda potência de dois $n\ge n_0$ e todo inteiro
\begin{equation}
 H\ge\max\{8n^{\nu_s},32(R+1)\},
 \label{lp:eq:height-threshold}
\end{equation}
há $g_{n,H}\in\Aclass_{n,H}([0,1])$ com
\begin{equation}
 \|f-g_{n,H}\|_p\le9C_W\,\wt{f}{n^{-1}}{p}+\frac{16(R+1)}{H}.
 \label{lp:eq:jackson-p}
\end{equation}
\end{teorema}

A prova constrói um quase-interpolante em $L^p$, obtém controle dos seus valores nodais e quantiza esses valores sob o orçamento $H$.

\subsection{Projeção afim local e quase-interpolação}

Para $I=[a,a+\ell]$, $\ell>0$, ponha $u_x=(x-a)/\ell$. Definimos
\[
 P_I h(x)=\int_I K_I(x,y)h(y)\dd y,
 \qquad
 K_I(x,y)=\frac{1+3(2u_x-1)(2u_y-1)}{\ell}.
\]

\begin{lema}[Projeção afim local]
\label{lp:lem:projection}
O operador $P_I$ reproduz funções afins e coincide em $L^2(I)$ com a projeção ortogonal sobre elas. Se $e$ é um extremo de $I$, então
\begin{equation}
 |(P_I h)(e)|\le4\ell^{-1/p}\|h\|_{L^p(I)},\qquad
 |(P_I h)(e)|\le2\|h\|_{L^\infty(I)}\quad(h\in L^\infty(I)).
 \label{lp:eq:projection-bounds}
\end{equation}
\end{lema}

\begin{proof}
As funções $\ell^{-1/2}$ e $\sqrt{3/\ell}(2u_x-1)$ formam uma base ortonormal do espaço das funções afins em $L^2(I)$. A soma dos produtos das funções da base em $x$ e $y$ dá $K_I(x,y)$, o que estabelece a projeção e a reprodução. Nos extremos, $\ell K_I(e,y)$ é $4-6u_y$ ou $6u_y-2$. Seu módulo é limitado por $4$, e Hölder dá a primeira estimativa. A integral do módulo de cada uma dessas funções em $u_y\in[0,1]$ vale $5/3$, o que dá a segunda.
\end{proof}

Para $n\ge2$, sejam $t_j=j/n$, $0\le j\le n$, e $I_k=[t_k,t_{k+1}]$, $0\le k<n$. Seja $\varphi_j$ a função chapéu contínua e afim em cada célula, de valor $1$ em $t_j$ e $0$ nos demais nós. Com $c(j)=\min\{j,n-1\}$, definimos
\begin{equation}
 \lambda_j(f)=(P_{I_{c(j)}}f)(t_j),
 \qquad Q_nf=\sum_{j=0}^n\lambda_j(f)\varphi_j .
 \label{lp:eq:quasi}
\end{equation}
Usaremos $J_k=I_k\cup I_{k+1}$ para $k<n-1$ e $J_{n-1}=I_{n-2}\cup I_{n-1}$. Escrevemos $E(f,J)_p=\inf_{a\,\mathrm{afim}}\|f-a\|_{L^p(J)}$.

\begin{lema}[Quase-interpolação por projeções locais]
\label{lp:lem:quasi}
Para $f\in L^p$ valem as seguintes propriedades.
\begin{enumerate}[label=(\alph*)]
 \item $Q_n$ é linear, produz funções contínuas e afins em cada célula $I_k$ e reproduz funções afins.
 \item $\|f-Q_nf\|_{L^p(I_k)}\le9E(f,J_k)_p$.
 \item Se $f$ é limitada, $|\lambda_j(f)|\le2\|f\|_\infty$.
 \item Se $f$ é contínua, $|\lambda_{j+1}(f)-\lambda_j(f)|\le5\omega_1(f,n^{-1})_\infty$ para $0\le j<n$.
\end{enumerate}
\end{lema}

\begin{proof}
A fórmula \eqref{lp:eq:quasi} dá a linearidade e a estrutura por células. Para $a$ afim, $\lambda_j(a)=a(t_j)$ pelo \cref{lp:lem:projection}, logo $Q_na=a$. Em $I_k$, os únicos termos de $Q_n$ são os associados a $t_k$ e $t_{k+1}$. Para toda função afim $a$, a primeira estimativa de \eqref{lp:eq:projection-bounds} e $\|\varphi_j\|_{L^p(I_k)}\le n^{-1/p}$ dão
\[
 \|Q_n(f-a)\|_{L^p(I_k)}
 \le4\|f-a\|_{L^p(I_k)}+4\|f-a\|_{L^p(I_{c(k+1)})}.
\]
Como $f-Q_nf=(f-a)-Q_n(f-a)$ e ambas as células estão em $J_k$, segue a cota $9\|f-a\|_{L^p(J_k)}$. O ínfimo em $a$ dá (b).

A segunda estimativa de \eqref{lp:eq:projection-bounds} dá (c). Aplicada a $f-f(t_j)$ em $I_{c(j)}$, ela fornece
$|\lambda_j(f)-f(t_j)|\le2\omega_1(f,n^{-1})_\infty$. Somando as estimativas nos dois nós e a de $|f(t_{j+1})-f(t_j)|$, obtemos (d).
\end{proof}

\subsection{Desigualdade de Whitney em partições}

A estimativa a seguir controla a soma dos erros locais por um único módulo global. O controle do erro de aproximação polinomial por diferenças finitas remonta a \citet{artigo_2-whitney1957} e integra a teoria da aproximação por splines \citep{artigo_2-devorelorentz1993}. A demonstração a seguir explicita uma constante que depende apenas de $p$.

\begin{lema}[Whitney de ordem dois em partições]
\label{lp:lem:whitney}
Existe $C_W=C_W(p)\ge1$ tal que, para $n\ge2$ e $f\in L^p$,
\begin{equation}
 \left(\sum_{k=0}^{n-1}E(f,J_k)_p^p\right)^{1/p}
 \le C_W\,\wt{f}{n^{-1}}{p}.
 \label{lp:eq:whitney-partition}
\end{equation}
\end{lema}

\begin{proof}
Escreva $J_k=[a,a+8\delta]$, com $\delta=1/(4n)$, e tome
$S_k=[a+\delta,a+7\delta]$. Para $|\rho|\le\delta$ e $x\in S_k$, defina $W_\rho(x)=f(x+\rho)-2f(x)+f(x-\rho)$. Consideremos a média dupla
\[
 g(x)=\delta^{-2}\int_0^\delta\!\!\int_0^\delta
            f(x-\delta+u+v)\dd v\dd u.
\]
Se $U(x)=\int_a^x f(y)\dd y$ e $V(x)=\int_a^x U(y)\dd y$, então
$g(x)=\delta^{-2}\big(V(x+\delta)-2V(x)+V(x-\delta)\big)$ em $S_k$.
Portanto $g'$ é absolutamente contínua e $g''=\delta^{-2}W_\delta$ quase em toda parte nesse intervalo. Para a reta tangente $A$ de $g$ em $a+4\delta$, a fórmula de Taylor com resto integral, seguida de Hölder, dá
\[
 \|g-A\|_{L^p(S_k)}
 \le (6\delta)^2\|g''\|_{L^p(S_k)}
 =36\|W_\delta\|_{L^p(S_k)} .
\]

A reflexão $(u,v)\mapsto(\delta-u,\delta-v)$ na integral que define $g$ dá
\[
 2\delta^2(g(x)-f(x))
 =\int_0^\delta\!\!\int_0^\delta W_{u+v-\delta}(x)\dd v\dd u .
\]
Aplicando Jensen à média no quadrado e Tonelli, obtemos
\[
 \|f-g\|_{L^p(S_k)}^p
 \le \frac{1}{2^p\delta^2}\int_0^\delta\!\!\int_0^\delta
        \|W_{u+v-\delta}\|_{L^p(S_k)}^p\dd v\dd u .
\]

Ponha $z=f-A$. Como a segunda diferença de $A$ se anula, na faixa $[a,a+\delta]$ vale
\[
 z(x)=2z(x+\delta)-z(x+2\delta)+W_\delta(x+\delta).
\]
Na faixa $[a+7\delta,a+8\delta]$ vale a identidade com $x-\delta$ e $x-2\delta$ no lugar de $x+\delta$ e $x+2\delta$. Em ambas as faixas, os argumentos do lado direito pertencem a $S_k$. Pela convexidade de $t\mapsto t^p$ em $[0,\infty)$,
$(\sum_{i=1}^j b_i)^p\le j^{p-1}\sum_{i=1}^j b_i^p$ para $b_i\ge0$. Aplicamos a desigualdade com três termos em cada faixa e majoramos $3^{p-1}$ por $4^{p-1}$. Após as translações para $S_k$, cada faixa contribui com um fator $2^p+1$ na integral de $|z|^p$. Assim, com $c_4=4^{p-1}$ e $M_p=1+2c_4(2^p+1)$,
\[
 \|z\|_{L^p(J_k)}^p
 \le M_p\|z\|_{L^p(S_k)}^p+2c_4\|W_\delta\|_{L^p(S_k)}^p.
\]
Usando a desigualdade de convexidade com $j=2$ em $z=(f-g)+(g-A)$, concluímos que
\begin{multline*}
 E(f,J_k)_p^p\le
 K_1\|W_\delta\|_{L^p(S_k)}^p\\
 {}+\frac{K_2}{2^p\delta^2}\int_0^\delta\!\!\int_0^\delta
 \|W_{u+v-\delta}\|_{L^p(S_k)}^p\dd v\dd u ,
\end{multline*}
em que $K_2=M_p2^{p-1}$ e $K_1=36^pK_2+2c_4$.

Cada ponto pertence a no máximo quatro intervalos $J_k$. Para $0<|\rho|\le\delta$, temos $S_k\subseteq[|\rho|,1-|\rho|]$ e $W_\rho(x)=\Delta_{|\rho|}^2f(x-|\rho|)$. Logo
\[
 \sum_k\|W_\rho\|_{L^p(S_k)}^p
 \le4\|\Delta_{|\rho|}^2f\|_{L^p([0,1-2|\rho|])}^p
 \le4\wt{f}{n^{-1}}{p}^p .
\]
Para $\rho=0$ a soma é nula. Somando as estimativas locais e usando a área $\delta^2$ do quadrado de integração, obtemos \eqref{lp:eq:whitney-partition} com
$C_W=(4K_1+4K_2)^{1/p}$.
\end{proof}

\begin{corolario}[Erro global de quase-interpolação]
\label{lp:cor:quasi-error}
Para todo $f\in L^p$ e todo $n\ge2$,
\begin{equation}
 \|f-Q_nf\|_p\le9C_W\,\wt{f}{n^{-1}}{p}.
 \label{lp:eq:quasi-error}
\end{equation}
\end{corolario}

\begin{proof}
Somamos as potências $p$ das estimativas locais do \cref{lp:lem:quasi}, uma vez que as células $I_k$ são disjuntas a menos dos extremos, e aplicamos \eqref{lp:eq:whitney-partition}.
\end{proof}

\subsection{Imersão contínua e controle da inclinação}

\begin{lema}[Imersão e módulo de primeira ordem]
\label{lp:lem:embedding}
Sejam $1/p<s<2$, $\alpha=s-1/p$ e $f\in\Bes{s}{p}{\infty}$. A classe $f$ tem um único representante contínuo $F$, com
$\|F\|_\infty\le C\|f\|^*_{\Bes{s}{p}{\infty}}$. Para $0<t\le1/2$, vale
\begin{equation}
 \omega_1(F,t)_\infty
 \le C\|f\|^*_{\Bes{s}{p}{\infty}}
 \begin{cases}
 t^\alpha,&0<\alpha<1,\\
 t\log(e/t),&\alpha=1,\\
 t,&\alpha>1.
 \end{cases}
 \label{lp:eq:omega1-regimes}
\end{equation}
A constante $C$ depende apenas de $p$ e $s$.
\end{lema}

\begin{proof}
Para uma função $v$ contínua e afim em cada célula de uma malha uniforme de passo $m^{-1}$, a equivalência de normas no espaço das funções afins de uma célula, seguida da mudança de escala, dá
\begin{equation}
 \|v\|_\infty\le C_p m^{1/p}\|v\|_p,\qquad \|v'\|_\infty\le2m\|v\|_\infty,
 \label{lp:eq:nikolskii}
\end{equation}
pois a inclinação em cada célula é $m$ vezes a diferença de valores nos extremos, com a derivada tomada no interior das células.

Ponha $S=\|f\|^*_{\Bes{s}{p}{\infty}}$ e $v_k=Q_{2^{k+1}}f-Q_{2^k}f$ para $k\ge1$. Cada $v_k$ é afim por células na malha de passo $2^{-(k+1)}$. Por \eqref{lp:eq:quasi-error}, pela monotonicidade de $t\mapsto\wt{f}{t}{p}$ e por \eqref{lp:eq:besov-p},
\[
 \|v_k\|_p\le18C_W\,\wt{f}{2^{-k}}{p}\le18C_W 2^{-ks}S.
\]
Assim $\|v_k\|_\infty\le C2^{-k\alpha}S$ e $\|v_k'\|_\infty\le C2^{k(1-\alpha)}S$. A série $Q_2f+\sum_{k\ge1}v_k$ converge uniformemente para uma função contínua $F$. A estimativa \eqref{lp:eq:quasi-error} também dá $Q_{2^r}f\to f$ em $L^p$, logo $F=f$ quase em toda parte. Dois representantes contínuos da mesma classe coincidem, pois todo subintervalo aberto não vazio tem medida positiva. A cota de $\|F\|_\infty$ resulta da série geométrica e de $\|Q_2f\|_\infty\le8\|f\|_p$, por \eqref{lp:eq:projection-bounds}.

Para $r\ge1$, a cauda da série dá $\|F-Q_{2^r}f\|_\infty\le C2^{-r\alpha}S$. Somando as cotas de inclinação,
\[
 \omega_1(F,2^{-r})_\infty
 \le 2\|F-Q_{2^r}f\|_\infty+2^{-r}\|(Q_{2^r}f)'\|_\infty
 \le CS\left(2^{-r\alpha}+2^{-r}\Bigl(1+\sum_{k=1}^{r-1}2^{k(1-\alpha)}\Bigr)\right).
\]
O fator entre parênteses é limitado por um múltiplo de $2^{-r\alpha}$, de $r2^{-r}$ ou de $2^{-r}$, conforme $\alpha<1$, $\alpha=1$ ou $\alpha>1$. Para $t$ entre duas escalas diádicas consecutivas, a monotonicidade de $\omega_1$ dá \eqref{lp:eq:omega1-regimes}.
\end{proof}

\subsection{Quantização nodal e prova da estimativa de Jackson}

\begin{lema}[Interpolação racional com orçamento de altura]
\label{lp:prop:jackson-proof}
Seja $n=2^r\ge2$ e sejam $y_0,\ldots,y_n$ dados reais com $|y_j|\le M$ e $|y_{j+1}-y_j|\le1$. Se $H\ge\max\{8n,16(M+1)\}$, o interpolante afim por células $u$ dos dados $(j/n,y_j)$ admite um aproximante $g\in\Aclass_{n,H}([0,1])$ tal que
\[
 \|u-g\|_\infty\le\frac{8(M+1)}{H}.
\]
\end{lema}

\begin{proof}
Seja $Q$ a maior potência de dois com $8(M+1)Q\le H$. Então $Q\ge1$ e $H<16(M+1)Q$. Para cada $j$, escolha o ponto $v_j$ mais próximo de $y_j$ na grade
$Q^{-1}(\Z+\tfrac12(j\bmod 2))$, com desempate pelo menor valor. Temos $|v_j-y_j|\le1/(2Q)$. Escrevendo $v_j=b_j/(2Q)$, temos $b_j\in\Z$, $b_j\equiv j\pmod 2$ e $|b_j|\le2QM+1$.

O interpolante $g$ dos valores $v_j$ é contínuo, e cada uma de suas $n$ células é uma folha da árvore diádica completa de profundidade $r$. Na folha de índice $j$, o valor central e a inclinação são
\[
 \frac{b_j+b_{j+1}}{4Q},
 \qquad
 \frac{n(b_{j+1}-b_j)}{2Q},
\]
com inclinação não nula pela paridade oposta dos numeradores. Para o valor central, o denominador e o módulo do numerador são limitados por $4Q$ e $4QM+2$, respectivamente, ambos no máximo $H$. Se $n\le2Q$, a inclinação pode ser escrita com numerador $b_{j+1}-b_j$ e denominador $2Q/n$, inteiros porque $n$ e $Q$ são potências de dois. As duas alturas ficam limitadas por $H$, visto que $|b_{j+1}-b_j|\le4QM+2$ e $2Q/n\le2Q$. Se $n>2Q$, a inclinação é inteira e seu módulo é no máximo
\[
 n\bigl(|y_{j+1}-y_j|+1/Q\bigr)\le2n\le H.
\]
A redução das frações preserva essas cotas. Como $r+1\le n$ e $j+1\le n$, a altura de cada folha é no máximo $H$. Por fim, $u-g$ é afim em cada célula e seu módulo é limitado pelo maior erro nos nós. Portanto $\|u-g\|_\infty\le1/(2Q)\le8(M+1)/H$.
\end{proof}

\begin{proof}[Demonstração do \cref{lp:hyp:jackson}]
Seja $F$ o representante contínuo de $f$ do \cref{lp:lem:embedding}. Os valores $y_j=\lambda_j(F)$ satisfazem $|y_j|\le2R$ pelo \cref{lp:lem:quasi}. O mesmo lema dá
$|y_{j+1}-y_j|\le5\omega_1(F,n^{-1})_\infty$. Pela continuidade uniforme de $F$, existe $n_0\ge2$ tal que essa cota é no máximo $1$ para $n\ge n_0$.

A condição \eqref{lp:eq:height-threshold} implica $H\ge8n$ e $H\ge16(2R+1)$. Aplicamos o \cref{lp:prop:jackson-proof} com $M=2R$ e $u=Q_nF$. Como $F=f$ quase em toda parte, $Q_nF=Q_nf$ e $\wt{F}{n^{-1}}{p}=\wt{f}{n^{-1}}{p}$. Pela desigualdade triangular, por \eqref{lp:eq:quasi-error} e pela normalização da medida de $[0,1]$,
\begin{align*}
 \|f-g_{n,H}\|_p
 &\le\|f-Q_nf\|_p+\|Q_nf-g_{n,H}\|_\infty\\
 &\le9C_W\,\wt{f}{n^{-1}}{p}+\frac{8(2R+1)}{H}
 \le9C_W\,\wt{f}{n^{-1}}{p}+\frac{16(R+1)}{H}.\qedhere
\end{align*}
\end{proof}

\begin{lema}[Quantização com salto limitado]
\label{lp:lem:nodal-jump}
Sejam $n=2^r\ge2$, $M\ge0$ e $J\ge0$, e sejam $y_0,\dots,y_n\in\mathbb R$ com $|y_j|\le M$ e $|y_{j+1}-y_j|\le J$ para $0\le j<n$. Se $H\in\mathbb N$ satisfaz
\[
 8n\le H,\qquad 16(M+1)\le H\qquad\text{e}\qquad n(J+1)\le H,
\]
então existe $g\in\Aclass_{n,H}([0,1])$ tal que $\|u-g\|_\infty\le 8(M+1)/H$, em que $u$ é o interpolante afim dos dados $(j/n,\,y_j)$.
\end{lema}

\begin{proof}
Seja $x=H/(16(M+1))\ge1$ e seja $Q$ a potência de dois com $x<Q\le2x$. Então $1/(2Q)\le8(M+1)/H$ e $8(M+1)Q\le H$. Distinguimos dois casos, conforme $n\le Q$ ou $n>Q$.

\emph{Caso $n\le Q$.} Como $H\ge8(M+1)Q$, o \cref{lp:lem:nodal-general} aplica-se com esse $Q$ e dá $g\in\Aclass_{n,H}([0,1])$ com $\|u-g\|_\infty\le1/(2Q)\le8(M+1)/H$.

\emph{Caso $n>Q$.} Como $n$ e $Q$ são potências de dois, $n\ge2Q$ e $n/(2Q)$ é inteiro. Escolha $b_j\in\mathbb Z$ de paridade $j\bmod2$ mais próximo de $2Qy_j$, com desempate pelo menor inteiro, e ponha $v_j=b_j/(2Q)$. Então $|v_j-y_j|\le1/(2Q)$ e $|b_j|\le2QM+1$. Seja $g$ o interpolante afim por células dos dados $(j/n,v_j)$ na árvore diádica completa de profundidade $r$. Na célula $j$, o valor central é $(b_j+b_{j+1})/(4Q)$ e a inclinação é
\begin{gather*}
 \sigma_j=n(v_{j+1}-v_j)=\frac{n}{2Q}\,(b_{j+1}-b_j)\in\mathbb Z\setminus\{0\},\\
 |\sigma_j|\le n\bigl(|y_{j+1}-y_j|+1/Q\bigr)\le n(J+1)\le H.
\end{gather*}
A paridade oposta de $b_j$ e $b_{j+1}$ garante a inclinação não nula. No valor central, o numerador tem módulo no máximo $4QM+2$ e o denominador é $4Q$, ambos limitados por $8(M+1)Q\le H$. A inclinação tem altura $|\sigma_j|\le H$. Os índices das folhas satisfazem $2\le r+1\le n\le H$ e $j+1\le n$. Logo $g\in\Aclass_{n,H}([0,1])$. Como $u$ e $g$ são afins em cada célula, $\|u-g\|_\infty\le\max_j|v_j-y_j|\le1/(2Q)\le8(M+1)/H$.
\end{proof}

\subsection{Quantização sem continuidade e Jackson para \texorpdfstring{$s\le1/p$}{s ≤ 1/p}}

Nesta subseção, $\theta_q=\max\{0,1-1/q\}$, com $\theta_\infty=1$. A construção para $s\le1/p$ controla os valores nodais $\lambda_j(f)$ a partir de projeções locais, inclusive nos regimes sem imersão contínua. A quantização a seguir exige $Q\ge n$ e usa a cota dos valores nodais para controlar a altura.

\begin{lema}[Interpolação racional sem controle dos saltos]
\label{lp:lem:nodal-general}
Seja $n=2^r\ge2$, seja $Q\ge n$ uma potência de dois e sejam $y_0,\ldots,y_n$ reais com $|y_j|\le M$. Se $H$ é um inteiro com $H\ge8(M+1)Q$, o interpolante afim por células $u$ dos dados $(j/n,y_j)$ admite $g\in\Aclass_{n,H}([0,1])$ tal que $\|u-g\|_\infty\le1/(2Q)$.
\end{lema}

\begin{proof}
Para cada $j$, escolha o ponto $v_j$ mais próximo de $y_j$ na grade $Q^{-1}(\Z+\tfrac12(j\bmod2))$, com desempate pelo menor valor. Então $|v_j-y_j|\le1/(2Q)$, e $v_j=b_j/(2Q)$ com $b_j\in\Z$, $b_j\equiv j\pmod2$ e $|b_j|\le2QM+1$. O interpolante $g$ dos valores $v_j$ é contínuo, e cada célula $[j/n,(j+1)/n]$ é uma folha da árvore diádica completa de profundidade $r$. Nessa folha, o valor central é $(b_j+b_{j+1})/(4Q)$ e a inclinação é $(b_{j+1}-b_j)/(2Q/n)$, não nula pela paridade oposta dos numeradores; aqui $2Q/n$ é inteiro porque $n\le Q$ são potências de dois. Os numeradores têm módulo no máximo $4QM+2$ e os denominadores são no máximo $4Q$. Como $H\ge8(M+1)Q$, todos esses números são no máximo $H$, e a redução das frações preserva as cotas. Além disso, $2\le r+1\le n\le Q\le H$ e $j+1\le n$, o que limita a altura de cada folha. Por fim, $u-g$ é afim em cada célula e seu módulo é limitado pelo maior erro nodal, $1/(2Q)$.
\end{proof}

\begin{lema}[Cota em sup dos quase-interpolantes]
\label{lp:lem:sup-bound}
Sejam $1\le p<\infty$, $0<s\le1/p$, $0<q\le\infty$ e $\delta=1/p-s$. Para $n=2^r\ge2$, ponha
\[
 \Lambda_n=n^{\delta}\quad(s<1/p),\qquad \Lambda_n=r^{\theta_q}\quad(s=1/p).
\]
Existe $C_0=C_0(p,s,q)\ge1$ tal que, para todo $f\in\Bes{s}{p}{q}$, com $S=\|f\|^*_{\Bes{s}{p}{q}}$, e todo $n=2^r\ge2$,
\[
 \max_{0\le j\le n}|\lambda_j(f)|=\|Q_nf\|_\infty\le C_0\,S\,\Lambda_n .
\]
\end{lema}

\begin{proof}
Ponha $a_k=2^{ks}\wt{f}{2^{-k}}{p}$. Por \eqref{lp:eq:besov-p}, $\|(a_k)_{k\ge0}\|_{\ell^q}\le S$ e $\|f\|_p\le S$. Para $k\ge1$, seja $v_k=Q_{2^{k+1}}f-Q_{2^k}f$. Como na demonstração do \cref{lp:lem:embedding}, $v_k$ é contínua e afim por células na malha de passo $2^{-(k+1)}$, e \eqref{lp:eq:quasi-error} com a monotonicidade do módulo dá $\|v_k\|_p\le18C_W\,\wt{f}{2^{-k}}{p}=18C_W2^{-ks}a_k$. Por \eqref{lp:eq:nikolskii} com $m=2^{k+1}$,
\[
 \|v_k\|_\infty\le C_E\,2^{k\delta}a_k,\qquad C_E=18\,C_W\,C_p\,2^{1/p}.
\]
Como $Q_nf=Q_2f+\sum_{k=1}^{r-1}v_k$ e $\|Q_2f\|_\infty\le8\|f\|_p\le8S$ por \eqref{lp:eq:projection-bounds},
\[
 \|Q_nf\|_\infty\le8S+C_E\sum_{k=1}^{r-1}2^{k\delta}a_k .
\]
Se $s<1/p$, usamos $a_k\le\|a\|_{\ell^\infty}\le\|a\|_{\ell^q}\le S$ e $\sum_{k=1}^{r-1}2^{k\delta}\le n^\delta/(2^\delta-1)$, o que dá a cota com $C_0=8+C_E/(2^\delta-1)$. Se $s=1/p$, então $\delta=0$ e a soma é $\sum_{k=1}^{r-1}a_k$. Para $q\le1$, vale $\sum_ka_k\le\|a\|_{\ell^1}\le\|a\|_{\ell^q}\le S$. Para $1<q\le\infty$, a desigualdade de Hölder dá $\sum_{k=1}^{r-1}a_k\le r^{1-1/q}\|a\|_{\ell^q}\le r^{\theta_q}S$. Em ambos os casos, como $r\ge1$, a cota vale com $C_0=8+C_E$.
\end{proof}

\begin{teorema}[Jackson para $s\le1/p$]
\label{lp:thm:jackson-low}
Sejam $1\le p<\infty$, $0<s\le1/p$ e $0<q\le\infty$, e sejam $\Lambda_n$ e $C_0$ como no \cref{lp:lem:sup-bound}. Sejam $f\in\Bes{s}{p}{q}$, $S=\|f\|^*_{\Bes{s}{p}{q}}$, $n=2^r\ge2$, $Q\ge n$ uma potência de dois e $H$ um inteiro com
\begin{equation}
 H\ge8C_0\,(S+1)\,\Lambda_n\,Q.
 \label{lp:eq:height-low}
\end{equation}
Então há $g\in\Aclass_{n,H}([0,1])$ com
\begin{equation}
 \|f-g\|_p\le9C_W\,\wt{f}{n^{-1}}{p}+\frac1{2Q}.
 \label{lp:eq:jackson-low}
\end{equation}
\end{teorema}

\begin{proof}
Ponha $y_j=\lambda_j(f)$ e $M=\|Q_nf\|_\infty$. Pelo \cref{lp:lem:sup-bound}, $|y_j|\le M\le C_0S\Lambda_n$. Como $C_0\ge1$ e $\Lambda_n\ge1$, vale $8(M+1)Q\le8C_0(S+1)\Lambda_nQ\le H$. O \cref{lp:lem:nodal-general}, aplicado a $u=Q_nf$, dá $g\in\Aclass_{n,H}([0,1])$ com $\|Q_nf-g\|_\infty\le1/(2Q)$. A desigualdade triangular, \eqref{lp:eq:quasi-error} e a normalização da medida de $[0,1]$ dão \eqref{lp:eq:jackson-low}.
\end{proof}

\section{Demonstração dos teoremas de caracterização e estabilidade}
\label{lp:sec:proved}

\subsection{Estimativa inversa}

\begin{lema}[Propriedades do módulo]
\label{lp:lem:modulus}
Para $f,g\in L^p$ e $t>0$,
\[
 \wt{f+g}{t}{p}\le\wt{f}{t}{p}+\wt{g}{t}{p},
 \qquad \wt{g}{t}{p}\le4\|g\|_p.
\]
Além disso, $t\mapsto\wt{f}{t}{p}$ é não decrescente.
\end{lema}

\begin{proof}
A primeira desigualdade decorre da linearidade de $\Delta_h^2$ e de Minkowski para cada $h$ admissível. Para a segunda, cada translação $x\mapsto x+jh$, $j\in\{0,1,2\}$, leva $[0,1-2h]$ a um subintervalo de $[0,1]$. A soma dos módulos dos coeficientes da diferença é $4$. A monotonicidade segue da inclusão dos conjuntos sobre os quais se toma o supremo em \eqref{lp:eq:modulus-p}.
\end{proof}

\begin{proposicao}[Desigualdade inversa de realização]
\label{lp:prop:inverse}
Para toda família $\mathcal G$ da \cref{lp:def:realization}, todo $f\in L^p$ e $N\ge2$,
\begin{equation}
 \wt{f}{N^{-1}}{p}\le4\Rfun{\mathcal G}{N}(f)_p.
 \label{lp:eq:inverse}
\end{equation}
Se $0\in\mathcal G_N$, vale também $\Rfun{\mathcal G}{N}(f)_p\le\|f\|_p$.
\end{proposicao}

\begin{proof}
Para $g\in\mathcal G_N$, o \cref{lp:lem:modulus} dá
\[
 \wt{f}{N^{-1}}{p}
 \le4\|f-g\|_p+\wt{g}{N^{-1}}{p}
 \le4\bigl(\|f-g\|_p+\wt{g}{N^{-1}}{p}\bigr).
\]
O ínfimo em $g$ dá \eqref{lp:eq:inverse}, e a escolha de $g=0$ dá a cota superior.
\end{proof}

\begin{corolario}[Inclusão em Besov]
\label{lp:cor:inclusion}
Para toda família $\mathcal G$ da \cref{lp:def:realization}, $0<s<2$ e $0<q\le\infty$, vale
$\Rstr{s}{q}{\mathcal G}_p\subseteq\Bes{s}{p}{q}$.
\end{corolario}

\begin{proof}
Multiplicamos \eqref{lp:eq:inverse} por $2^{rs}$ para $N=2^r$, $r\ge1$. O termo $r=0$ da sequência em \eqref{lp:eq:besov-p} é finito, pois $\wt{f}{1}{p}\le4\|f\|_p$.
\end{proof}

\subsection{Estimativa direta e agendas de altura}

\begin{proposicao}[Critério de somabilidade da agenda]
\label{lp:thm:besov-conditional}
Sejam $1/p<s<2$ e $0<q\le\infty$. Se $H(N)\ge8N^{\nu_s}$ para $N\ge1$ e
\begin{equation}
 \left(\frac{2^{rs}}{H(2^r)}\right)_{r\ge1}\in\ell^q,
 \label{lp:eq:summability}
\end{equation}
então $\Bes{s}{p}{q}=\Rstr{s}{q}{\Library,H}_p$.
\end{proposicao}

\begin{proof}
A inclusão em Besov é o \cref{lp:cor:inclusion}. Para a inclusão direta, tome $f\in\Bes{s}{p}{q}\subseteq\Bes{s}{p}{\infty}$ e ponha $R=\|f\|_\infty$ no representante contínuo. Para $N=2^r$ suficientemente grande, $H(N)$ satisfaz \eqref{lp:eq:height-threshold} e o \cref{lp:hyp:jackson} dá $g=g_{N,H(N)}\in\Sig{N}$. Pelo \cref{lp:lem:modulus},
$\wt{g}{N^{-1}}{p}\le\wt{f}{N^{-1}}{p}+4\|f-g\|_p$. Assim
\begin{align}
 \Rfun{\Library,H}{N}(f)_p
 &\le \wt{f}{N^{-1}}{p}+5\|f-g\|_p\notag\\
 &\le (1+45 C_W)\wt{f}{N^{-1}}{p}+\frac{80(R+1)}{H(N)}.
 \label{lp:eq:direct-bound}
\end{align}
Após a multiplicação por $N^s$, o primeiro termo pertence a $\ell^q$ por \eqref{lp:eq:besov-p}, e o segundo por \eqref{lp:eq:summability}. Para $0<q<\infty$, a soma permanece em $\ell^q$ pela desigualdade $(a+b)^q\le2^q(a^q+b^q)$, $a,b\ge0$. Para $q=\infty$, basta tomar o supremo. Nos orçamentos diádicos restantes, a escolha do aproximante nulo dá $\Rfun{\Library,H}{N}(f)_p\le\|f\|_p$.
\end{proof}

\begin{proof}[Demonstração do \cref{lp:conj:besov}, item (i)]
Ambas as agendas satisfazem $H(N)\ge8N^{\nu_s}$. Para $q=\infty$ e $H_{s,0}$, temos $N^s/H_{s,0}(N)\le1/8$. Se $q<\infty$ e $s<1$, obtemos
$N^s/H_{s,0}(N)\le2^{-r(1-s)}/8$ para $N=2^r$, o que dá uma sequência em $\ell^q$.
Quando $q<\infty$ e $s\ge1$, vale
\[
 \frac{2^{rs}}{H_{s,\kappa}(2^r)}\le\frac{1}{8}(1+r\log 2)^{-\kappa}.
\]
A série das potências $q$ dessa cota converge para $\kappa q>1$. Por fim,
$2^{rs}/\widetilde H_{s,\eta}(2^r)\le2^{-r\eta}/8$ para $\eta>0$, uma vez que $\max\{1,s+\eta\}\ge s+\eta$. Em cada caso, aplica-se a \cref{lp:thm:besov-conditional}.
\end{proof}

\subsection{Transferência para famílias aninhadas}

\begin{proposicao}[Transferência da estimativa direta]
\label{lp:thm:stability-conditional}
Sejam $1/p<s<2$ e $0<q\le\infty$. Suponha que $H(N)\ge8N^{\nu_s}$ para $N\ge1$ e que valha \eqref{lp:eq:summability}. Se $\mathcal N$ satisfaz \eqref{lp:eq:nesting} e \eqref{lp:eq:contains}, então $\Bes{s}{p}{q}\subseteq\Rstr{s}{q}{\mathcal N}_p$.
\end{proposicao}

\begin{proof}
Fixe $f\in\Bes{s}{p}{q}$ e ponha $R=\|f\|_\infty$ no representante contínuo. Seja $e=\lceil\log_2 c\rceil$ e $d=2^e\ge c$. Para $M=2^r$ suficientemente grande, tome $n=M/d$. Como $\lceil cn\rceil\le M$, temos
$\Sig{n}\subseteq\mathcal N_{\lceil cn\rceil}\subseteq\mathcal N_M$.
O aproximante de Jackson $g=g_{n,H(n)}$ pertence, portanto, a $\mathcal N_M$.
A monotonicidade em escala dá $\wt{g}{M^{-1}}{p}\le\wt{g}{n^{-1}}{p}$. Aplicando a estimativa dos termos em \eqref{lp:eq:direct-bound} na escala $n$, obtemos
\[
 M^s\Rfun{\mathcal N}{M}(f)_p
 \le d^s\left((1+45C_W)n^s\wt{f}{n^{-1}}{p}+\frac{80(R+1)n^s}{H(n)}\right).
\]
A passagem de $M$ a $n$ desloca o índice diádico de $r$ para $r-e$. Os dois termos entre parênteses pertencem a $\ell^q$ por \eqref{lp:eq:besov-p} e \eqref{lp:eq:summability}. Nos orçamentos restantes, $\Rfun{\mathcal N}{M}(f)_p\le\|f\|_p$ porque $0\in\mathcal N_M$.
\end{proof}

\begin{proof}[Demonstração do \cref{lp:conj:stability}]
Se $s>1/p$, as agendas do \cref{lp:conj:besov} satisfazem o critério \eqref{lp:eq:summability}, como se mostrou na subseção anterior, e a \cref{lp:thm:stability-conditional} dá a inclusão direta. Se $s\le1/p$, a inclusão direta é a \cref{lp:thm:low-transfer}. Em ambos os casos, o \cref{lp:cor:inclusion}, aplicado a $\mathcal N$, dá a inclusão inversa.
\end{proof}

\begin{proposicao}[Transferência sob inclusão aproximada para $s>1/p$]
\label{lp:thm:approx-transfer}
Sejam $1/p<s<2$ e $0<q\le\infty$. Suponha que $H(N)\ge8N^{\nu_s}$ para $N\ge1$ e que valham \eqref{lp:eq:summability} e \eqref{lp:eq:approx-summable}. Se $\mathcal N$ satisfaz \eqref{lp:eq:nesting} e \eqref{lp:eq:approx-contains}, então $\Bes{s}{p}{q}\subseteq\Rstr{s}{q}{\mathcal N}_p$.
\end{proposicao}

\begin{proof}
Fixe $f\in\Bes{s}{p}{q}$ e ponha $R=\|f\|_\infty$ no representante contínuo. Tome $e=\lceil\log_2c\rceil$ e $d=2^e\ge c$. Para $M=2^r$ suficientemente grande e $n=M/d$, o \cref{lp:hyp:jackson} fornece $g=g_{n,H(n)}\in\Sig{n}$ satisfazendo \eqref{lp:eq:jackson-p} com $H=H(n)$. Por \eqref{lp:eq:approx-contains} e \eqref{lp:eq:nesting}, existe $\tilde g\in\mathcal N_{\lceil cn\rceil}\subseteq\mathcal N_M$ tal que $\|g-\tilde g\|_p\le\varepsilon_n$. Logo $\|f-\tilde g\|_p\le\|f-g\|_p+\varepsilon_n$. Como $n\le M$, a monotonicidade do módulo e o \cref{lp:lem:modulus} dão
\[
 \wt{\tilde g}{M^{-1}}{p}
 \le\wt{\tilde g}{n^{-1}}{p}
 \le\wt{f}{n^{-1}}{p}+4\|f-\tilde g\|_p
 \le\wt{f}{n^{-1}}{p}+4\|f-g\|_p+4\varepsilon_n.
\]
Somando com $\|f-\tilde g\|_p$ e multiplicando por $M^s=d^s n^s$, obtemos
\[
 M^s\Rfun{\mathcal N}{M}(f)_p
 \le d^s\left((1+45C_W)n^s\wt{f}{n^{-1}}{p}+\frac{80(R+1)n^s}{H(n)}+5\,n^s\varepsilon_n\right).
\]
Os três termos entre parênteses pertencem a $\ell^q$ por \eqref{lp:eq:besov-p}, \eqref{lp:eq:summability} e \eqref{lp:eq:approx-summable}, após deslocamento de $e$ no índice diádico. Para $0<q<\infty$, a soma finita permanece em $\ell^q$; para $q=\infty$, toma-se o supremo. Nos orçamentos restantes, $0\in\mathcal N_M$ dá $\Rfun{\mathcal N}{M}(f)_p\le\|f\|_p$.
\end{proof}

\subsection{Regimes \texorpdfstring{$s\le1/p$}{s ≤ 1/p}}

\begin{proposicao}[Transferência para $s\le1/p$]
\label{lp:thm:low-transfer}
Sejam $1\le p<\infty$, $0<s\le1/p$ e $0<q\le\infty$. Seja $H=\widehat H_s$ se $s<1/p$, e $H=\widehat H_{q,\kappa}$ se $s=1/p$, com $\kappa\ge0$ e, quando $p=1$ e $q<\infty$, $\kappa>1/q$. Se $\mathcal N$ satisfaz \eqref{lp:eq:nesting} e \eqref{lp:eq:contains} para algum $c\ge1$, então $\Bes{s}{p}{q}\subseteq\Rstr{s}{q}{\mathcal N}_p$.
\end{proposicao}

\begin{proof}
Fixe $f\in\Bes{s}{p}{q}$ e ponha $S=\|f\|^*_{\Bes{s}{p}{q}}$, $K=8C_0(S+1)$ e $\delta=1/p-s$. Se $s<1/p$, ponha $\kappa_0=0$. Se $s=1/p$, ponha $\kappa_0=\kappa$ e $\vartheta=\theta_q+\kappa$. Para $n=2^r\ge2$, seja $Q(n)$ a menor potência de dois com $Q(n)\ge nr^{\kappa_0}$. Então $n\le Q(n)$ e $nr^{\kappa_0}\le Q(n)\le2nr^{\kappa_0}$.

\emph{Escolha do deslocamento.} Existe um inteiro $e\ge0$ tal que $d=2^e$ satisfaz
\begin{equation}
 K\,\Lambda_n\,Q(n)\le H(dn)\qquad(n=2^r\ge2).
 \label{lp:eq:shift}
\end{equation}
De fato, se $s<1/p$, o lado esquerdo é $Kn^{1+\delta}$ e $H(dn)\ge8d^{1+\delta}n^{1+\delta}$, e basta $8d^{1+\delta}\ge K$. Se $s=1/p$, o lado esquerdo é no máximo $2Knr^{\vartheta}$. Como $\log(en)=1+r\log2\ge r\log2$ e $1/\log2<2$, vale $r^\vartheta\le2^\vartheta[\log(en)]^\vartheta$. Por outro lado, $H(dn)\ge8dn[\log(en)]^\vartheta$, e basta $8d\ge2^{\vartheta+1}K$.

\emph{Aproximante.} Seja $d_c=2^{\lceil\log_2c\rceil}\ge c$ e $D=d_cd$. Para $M=2^\mu\ge2D$, ponha $N'=M/d_c$ e $n=M/D=N'/d$, de modo que $n\ge2$ é potência de dois. Por \eqref{lp:eq:shift}, $H(N')\ge K\Lambda_nQ(n)$, e o \cref{lp:thm:jackson-low}, com $Q=Q(n)$ e $H=H(N')$, dá
\[
 g\in\Aclass_{n,H(N')}([0,1])\subseteq\Aclass_{N',H(N')}([0,1])\subseteq\Sig{N'},
 \qquad
 \|f-g\|_p\le9C_W\,\wt{f}{n^{-1}}{p}+\frac1{2Q(n)} .
\]
Como $\lceil cN'\rceil\le d_cN'=M$, \eqref{lp:eq:contains} e \eqref{lp:eq:nesting} dão $g\in\mathcal N_{\lceil cN'\rceil}\subseteq\mathcal N_M$. Pelo \cref{lp:lem:modulus} e pela monotonicidade do módulo, $\wt{g}{M^{-1}}{p}\le\wt{f}{n^{-1}}{p}+4\|f-g\|_p$. Logo
\[
 M^s\Rfun{\mathcal N}{M}(f)_p
 \le D^s\Bigl((1+45C_W)\,n^s\wt{f}{n^{-1}}{p}+\frac52\,\frac{n^s}{Q(n)}\Bigr).
\]

\emph{Somabilidade.} A passagem de $M$ a $n$ desloca o índice diádico em $\log_2D$. O primeiro termo pertence a $\ell^q$ por \eqref{lp:eq:besov-p}. Para o segundo, $n^s/Q(n)\le n^{s-1}r^{-\kappa_0}$ com $n=2^r$. Se $s<1$, a sequência $2^{-r(1-s)}$ pertence a $\ell^q$ para todo $q$, e $r^{-\kappa_0}\le1$; isso cobre $s<1/p$ e também $s=1/p$ com $p>1$. Resta $s=1/p=1$, em que a sequência é $r^{-\kappa}$: ela é limitada para $q=\infty$ e pertence a $\ell^q$ quando $\kappa q>1$. Para $0<q<\infty$, a soma de duas sequências de $\ell^q$ permanece em $\ell^q$. Nos orçamentos $M$ restantes, $\Rfun{\mathcal N}{M}(f)_p\le\|f\|_p$ porque $0\in\mathcal N_M$.
\end{proof}

\begin{proof}[Demonstração do \cref{lp:conj:besov}, itens (ii) e (iii)]
A inclusão $\Rstr{s}{q}{\Library,H}_p\subseteq\Bes{s}{p}{q}$ é o \cref{lp:cor:inclusion}. Para a inclusão inversa, a família $\mathcal N_N=\Sig{N}$, $N\ge1$, satisfaz \eqref{lp:eq:nesting}, porque as agendas $\widehat H_s$ e $\widehat H_{q,\kappa}$ são não decrescentes, e satisfaz \eqref{lp:eq:contains} com $c=1$. A \cref{lp:thm:low-transfer} dá $\Bes{s}{p}{q}\subseteq\Rstr{s}{q}{\mathcal N}_p=\Rstr{s}{q}{\Library,H}_p$.
\end{proof}

\begin{proposicao}[Transferência aproximada para $s\le1/p$]
\label{lp:thm:low-approx-transfer}
Sejam $p,s,q,H$ como na \cref{lp:thm:low-transfer}. Se $\mathcal N$ satisfaz \eqref{lp:eq:nesting} e \eqref{lp:eq:approx-contains} para algum $c\ge1$ e uma sequência $(\varepsilon_N)_{N\ge1}$ em $[0,\infty)$ que cumpre \eqref{lp:eq:approx-summable}, então $\Bes{s}{p}{q}\subseteq\Rstr{s}{q}{\mathcal N}_p$.
\end{proposicao}

\begin{proof}
Mantemos $Q(n)$, $d$, $d_c$, $D=d_cd$, $N'=M/d_c$ e $n=M/D$ como na prova da \cref{lp:thm:low-transfer}. Para $M=2^\mu\ge2D$, o aproximante $g\in\Sig{N'}$ ali obtido tem aproximante $\tilde g\in\mathcal N_{\lceil cN'\rceil}\subseteq\mathcal N_M$ com $\|g-\tilde g\|_p\le\varepsilon_{N'}$ por \eqref{lp:eq:approx-contains} e \eqref{lp:eq:nesting}. Usando $\|f-\tilde g\|_p\le\|f-g\|_p+\varepsilon_{N'}$, a monotonicidade do módulo ($n\le M$) e o \cref{lp:lem:modulus}, vem
\[
 \Rfun{\mathcal N}{M}(f)_p
 \le\wt{f}{n^{-1}}{p}+5\|f-g\|_p+5\varepsilon_{N'}.
\]
Como $M=Dn=d_cN'$, multiplicamos por $M^s$:
\[
 M^s\Rfun{\mathcal N}{M}(f)_p
 \le D^s\Bigl((1+45C_W)\,n^s\wt{f}{n^{-1}}{p}+\frac52\,\frac{n^s}{Q(n)}\Bigr)+5d_c^s\,(N')^s\varepsilon_{N'}.
\]
A primeira parcela pertence a $\ell^q$ pelo mesmo argumento da \cref{lp:thm:low-transfer}, enquanto a segunda está em $\ell^q$ por \eqref{lp:eq:approx-summable}, pois $N'=M/d_c$ apenas desloca o índice diádico em $\lceil\log_2c\rceil$. Nos orçamentos $M$ restantes, $0\in\mathcal N_M$ dá $\Rfun{\mathcal N}{M}(f)_p\le\|f\|_p$.
\end{proof}

\begin{proof}[Demonstração do \cref{lp:thm:approx-stability}]
Se $s>1/p$, o critério \eqref{lp:eq:summability} vale para as agendas do \cref{lp:conj:besov} e a \cref{lp:thm:approx-transfer} dá $\Bes{s}{p}{q}\subseteq\Rstr{s}{q}{\mathcal N}_p$. Se $s\le1/p$, essa inclusão é a \cref{lp:thm:low-approx-transfer}. Em ambos os casos, a inclusão inversa segue do \cref{lp:cor:inclusion}.
\end{proof}

\section{Cotas uniformes, complexidade em bits e agendas de log iterado}
\label{lp:sec:uniform}

Esta seção reúne quatro complementos formalizados em Lean (ver a \cref{lp:sec:lean}): cotas uniformes sobre bolas de Besov com uma única agenda (\cref{lp:sec:ball}), o piso aritmético e a complexidade em bits dos aproximantes (\cref{lp:sec:bits}) e a redução do fator logarítmico crítico por logaritmos iterados (\cref{lp:sec:loglog}). As demonstrações usam o \cref{lp:lem:nodal-general}, o \cref{lp:lem:sup-bound} e o critério de somabilidade da \cref{lp:thm:besov-conditional}. O ponto de partida é a versão em norma uniforme desses resultados na Parte~IV de \citet{lp-oliveiravol}, com a imersão contínua e a cota em sup dos quase-interpolantes substituindo ali a norma supremo.

\subsection{Cotas uniformes sobre bolas de Besov}
\label{lp:sec:ball}

Para $B\ge1$, escrevemos $\mathbb B^s_{p,q}(B)=\{f\in L^p:\|f\|^*_{\Bes{s}{p}{q}}\le B\}$. No \cref{lp:conj:besov}, a agenda é a mesma para todas as funções, e a prova pode usar limiares que dependem de $f$. Aqui, fixados $(p,s,q)$, as cotas valem desde $n=2$ para toda a bola, e a agenda inclui o fator $B+1$. Ponha $\mathsf L(N)=\log(eN)$ e $\theta_q=\max\{0,1-1/q\}$, com $\theta_\infty=1$.

\begin{definicao}[Dados de regime]
\label{lp:def:regime}
Sejam $1\le p<\infty$, $0<s<2$, $0<q\le\infty$ e $\kappa\ge0$. Para $n=2^r\ge2$, definimos os pares $(\nu',\kappa')$ e $(\hat\nu,\hat\kappa)$ e o fator $\Lambda_n$ por meio da \cref{lp:tab:regime}.
\begin{table}[H]
\caption{Parâmetros para cada regime de regularidade}
\label{lp:tab:regime}
\small\tablecolors
\begin{tabularx}{\linewidth}{l*{3}{>{\centering\arraybackslash}X}}
\toprule
\tablehead{Regime} & \tablehead{$(\nu\prime,\kappa\prime)$} & \tablehead{$(\hat\nu,\hat\kappa)$} & \tablehead{$\Lambda_n$} \\
\midrule
$s>1/p$ & $(\nu_s,\kappa)$ & $(\nu_s,\kappa)$ & $1$ \\
$s<1/p$ & $(1,0)$ & $(1+1/p-s,0)$ & $n^{1/p-s}$ \\
$s=1/p$ & $(1,\kappa)$ & $(1,\theta_q+\kappa)$ & $r^{\theta_q}$ \\
\bottomrule
\end{tabularx}\par
\tablecredit
\end{table}
Dizemos que $(p,s,q,\kappa)$ é \emph{admissível} quando $\kappa>1/q$ sempre que $q<\infty$ e $s\ge1$. Para $B\ge1$ e $A\ge1$, a \emph{agenda uniforme} é
\begin{equation}
 \widehat H^{[B]}_{A}(N)=\Bigl\lceil A\,(B+1)\,(N+1)^{\hat\nu}\,\mathsf L(N)^{\hat\kappa}\Bigr\rceil .
 \label{lp:eq:uniform-schedule}
\end{equation}
\end{definicao}

\begin{lema}[Altura e quantização uniformes]
\label{lp:lem:uniform-height}
Sejam $1\le p<\infty$, $0<s<2$, $0<q\le\infty$ e $\kappa\ge0$. Para $n=2^r\ge2$, seja $Q(n)$ a menor potência de dois com $Q(n)\ge n^{\nu'}\mathsf L(n)^{\kappa'}$. Existe $C_*\ge1$, dependendo só de $(p,s,q)$, tal que:
\begin{enumerate}[label=(\alph*)]
 \item $n\le Q(n)\le2n^{\nu'}\mathsf L(n)^{\kappa'}$, $\ 1/Q(n)\le n^{-s}$ e $\Lambda_nQ(n)\le4n^{\hat\nu}\mathsf L(n)^{\hat\kappa}$;
 \item se $f\in\mathbb B^s_{p,q}(B)$, então $\max_{0\le j\le n}|\lambda_j(f)|\le C_*B\Lambda_n$ para todo $n=2^r\ge2$;
 \item se $(p,s,q,\kappa)$ é admissível, então $\big(2^{rs}/Q(2^r)\big)_{r\ge1}\in\ell^q$, com norma limitada por uma constante $c_2$ que depende só de $(p,s,q,\kappa)$.
\end{enumerate}
\end{lema}

\begin{proof}
(a) Como $\nu'\ge1$ e $\mathsf L(n)\ge1$, vale $n^{\nu'}\mathsf L(n)^{\kappa'}\ge n$, e a menor potência de dois não inferior a $x\ge1$ é menor que $2x$. Além disso, $\nu'\ge s$: para $s>1/p$, $\nu_s\ge s$; para $s\le1/p$, $\nu'=1\ge s$. Logo $1/Q(n)\le n^{-s}$. Para a última cota, se $s>1/p$, então $\Lambda_nQ(n)\le2n^{\nu_s}\mathsf L(n)^\kappa$. Se $s<1/p$, então $Q(n)=n$ e $\Lambda_nQ(n)=n^{1+1/p-s}$. Se $s=1/p$, então $\mathsf L(n)=1+r\log2\ge r\log2$ e $1/\log2<2$ dão $r^{\theta_q}\le2^{\theta_q}\mathsf L(n)^{\theta_q}$, de modo que $\Lambda_nQ(n)\le2^{\theta_q+1}n\,\mathsf L(n)^{\theta_q+\kappa}\le4n\,\mathsf L(n)^{\theta_q+\kappa}$.

(b) Se $s>1/p$, o \cref{lp:lem:embedding} dá $\|F\|_\infty\le C_E\|f\|^*_{\Bes{s}{p}{\infty}}\le C_EB$ para o representante contínuo $F$, pois a norma em $\ell^\infty$ não excede a norma em $\ell^q$. O \cref{lp:lem:quasi}(c) dá $|\lambda_j(f)|=|\lambda_j(F)|\le2C_EB$. Se $s\le1/p$, o \cref{lp:lem:sup-bound} dá $|\lambda_j(f)|\le C_0B\Lambda_n$. Basta tomar $C_*=\max\{2C_E,C_0,1\}$.

(c) Por (a), $w_n=n^s/Q(n)\le n^{s-\nu'}\mathsf L(n)^{-\kappa'}$. Se $q=\infty$, então $w_n\le1$. Se $q<\infty$ e $s<1$, então $\nu'=1$ e $w_{2^r}\le2^{-r(1-s)}$, que pertence a $\ell^q$. Se $q<\infty$ e $s\ge1$, então $\kappa>1/q$ por admissibilidade, $\nu'=s$ e $\kappa'=\kappa$, de modo que $w_{2^r}\le(1+r\log2)^{-\kappa}$, que pertence a $\ell^q$ porque $\kappa q>1$.
\end{proof}

\begin{teorema}[Cotas uniformes sobre bolas]
\label{lp:thm:uniform}
Sejam $1\le p<\infty$, $0<s<2$, $0<q\le\infty$ e $\kappa\ge0$, e sejam $C_*$ e $c_2$ como no \cref{lp:lem:uniform-height}. Ponha $A=32(C_*+1)$ e $\widehat H=\widehat H^{[B]}_{A}$. Então, para todo $B\ge1$, todo $f\in\mathbb B^s_{p,q}(B)$ e todo $n=2^r\ge2$, existe $g\in\Aclass_{n,\widehat H(n)}([0,1])\subseteq\Sig{n}$, com agenda $\widehat H$, tal que
\begin{equation}
 \|f-g\|_p\le9C_W\,\wt{f}{n^{-1}}{p}+\frac1{2Q(n)}\le(9C_WB+1)\,n^{-s}.
 \label{lp:eq:uniform-approx}
\end{equation}
Se, além disso, $(p,s,q,\kappa)$ é admissível, existe $C=C(p,s,q,\kappa)$ tal que
\begin{equation}
 \sup_{f\in\mathbb B^s_{p,q}(B)}\Big\|\big(2^{rs}\Rfun{\Library,\widehat H}{2^r}(f)_p\big)_{r\ge1}\Big\|_{\ell^q}\le C\,(B+1).
 \label{lp:eq:uniform-rate}
\end{equation}
\end{teorema}

\begin{proof}
Fixe $f\in\mathbb B^s_{p,q}(B)$ e $n=2^r\ge2$. Ponha $M=C_*B\Lambda_n$ e $Q=Q(n)$. Pelo \cref{lp:lem:uniform-height}(b), $|\lambda_j(f)|\le M$ para $0\le j\le n$. Como $\Lambda_n\ge1$, o item (a) do mesmo lema dá
\[
 8(M+1)Q\le8(C_*B+1)\Lambda_nQ\le32(C_*+1)(B+1)\,n^{\hat\nu}\mathsf L(n)^{\hat\kappa}\le\widehat H(n).
\]
Como $Q\ge n$, o \cref{lp:lem:nodal-general} aplicado a $y_j=\lambda_j(f)$ dá $g\in\Aclass_{n,\widehat H(n)}([0,1])$ com $\|Q_nf-g\|_\infty\le1/(2Q)$. Pelo \cref{lp:cor:quasi-error}, $\|f-g\|_p\le9C_W\,\wt{f}{n^{-1}}{p}+1/(2Q)$. Por \eqref{lp:eq:besov-p}, $\wt{f}{n^{-1}}{p}\le Bn^{-s}$, e por (a), $1/(2Q)\le n^{-s}/2$. Isso prova \eqref{lp:eq:uniform-approx}.

Para \eqref{lp:eq:uniform-rate}, o \cref{lp:lem:modulus} dá $\wt{g}{n^{-1}}{p}\le\wt{f}{n^{-1}}{p}+4\|f-g\|_p$, de modo que, como na prova da \cref{lp:thm:besov-conditional},
\[
 \Rfun{\Library,\widehat H}{n}(f)_p\le(1+45C_W)\,\wt{f}{n^{-1}}{p}+\frac52\,\frac1{Q(n)}.
\]
Multiplicando por $n^s=2^{rs}$ e tomando a norma de $\ell^q$ em $r\ge1$, o primeiro termo contribui com no máximo $(1+45C_W)B$ por \eqref{lp:eq:besov-p}, e o segundo, com no máximo $\tfrac52c_2$ pelo \cref{lp:lem:uniform-height}(c). A desigualdade triangular em $\ell^q$ tem constante $2^{(1/q-1)_+}$, que não depende de $B$. Isso dá \eqref{lp:eq:uniform-rate}.
\end{proof}

\begin{observacao}
\label{lp:obs:uniform-q}
A taxa por escala \eqref{lp:eq:uniform-approx} vale com $\kappa=0$ para todo $q$. Na diagonal $s=1/p$, a agenda \eqref{lp:eq:uniform-schedule} ainda inclui $\mathsf L(N)^{\theta_q}$ para controlar os valores nodais. Quando $s\ge1$ e $q<\infty$, o fator adicional $\mathsf L(N)^\kappa$ com $\kappa>1/q$ torna a cota de quantização somável. Com $\kappa=0$, o peso $2^{rs}/Q(2^r)$ fica entre $1/2$ e $1$, portanto essa majoração não dá \eqref{lp:eq:uniform-rate}.
\end{observacao}

\begin{lema}[Saltos nodais uniformes]
\label{lp:lem:critical-jump}
Sejam $p>1$, $0<q\le\infty$ e $f$ com $\|f\|^*_{\Bes{1/p}{p}{q}}\le B$, $n=2^r\ge2$ e $\lambda_j=\lambda_j(f)$. Para $0\le j<n$,
\begin{equation}
 |\lambda_{j+1}-\lambda_j|\le\frac{4\,\sup_{[0,1]}|Q_2f|}{n}+\sum_{k=1}^{r-1}\frac{K_k}{n},
 \qquad K_k=8\,(2^{k+1})^{1/p+1}\,18\,C_W\,\wt{f}{2^{-k}}{p}.
 \label{lp:eq:critical-jump}
\end{equation}
Além disso, $K_k\le 144\,C_W\,2^{1/p}\,2^{k+1}\,B$ e $\sum_{k=1}^{r-1}K_k/n\le 288\,C_W\,2^{1/p}B$.
\end{lema}

\begin{proof}
Pela telescópica $Q_{2^r}f=Q_2f+\sum_{k=1}^{r-1}(Q_{2^{k+1}}f-Q_{2^k}f)$ e por $\lambda_j=Q_{2^r}f(j/n)$, a diferença $\lambda_{j+1}-\lambda_j$ é a soma das diferenças dos termos. Cada termo é afim em células de comprimento $1/2$ (para $Q_2f$) ou $2^{-(k+1)}$ (para as diferenças), e o intervalo $[j/n,(j+1)/n]$ está contido em uma dessas células. Aplicando \eqref{lp:eq:projection-bounds} a uma função afim na própria célula, pode-se tomar $C_p=4$ em \eqref{lp:eq:nikolskii}. Assim, para $m=2^{k+1}$, a inclinação de $v_k=Q_{2^{k+1}}f-Q_{2^k}f$ é limitada por $8m^{1/p+1}\|v_k\|_p\le K_k$, pela cota de quase-interpolação nos dois níveis. A inclinação de $Q_2f$ é limitada por $4\sup|Q_2f|$. A desigualdade $2^{k/p}\wt f{2^{-k}}p\le B$ dá a segunda parte.
\end{proof}

\begin{teorema}[Altura linear crítica para $s=1/p$ com $p>1$]
\label{lp:thm:critical-linear}
Sejam $p>1$, $s=1/p$ e $0<q<\infty$. Existem $A=A(p,C_W)\ge1$ e constantes $C_2,C_3>0$ que dependem só de $(p,q,C_W)$, tais que, para todo $B\ge1$, a agenda linear $H_B(N)=\lceil A(B+1)(N+1)\rceil$ satisfaz as seguintes propriedades.
\begin{enumerate}[label=(\alph*)]
 \item Para todo $f\in\mathbb B^{1/p}_{p,q}(B)$ e $n=2^r\ge2$, existe $g\in\Aclass_{n,H_B(n)}([0,1])$ com
 \[
  \|f-g\|_p\le 9C_W\,\wt{f}{n^{-1}}{p}+\frac{8(M_r+1)}{H_B(n)},\qquad M_r=\mathsf B_0(f)+144\,C_W\,2^{1/p}\,B\,r,
 \]
 em que $\mathsf B_0(f)=8\,2^{1/p}(\|f\|_p+9C_W\wt{f}{1/2}{p})$.
 \item Para todo $f\in\mathbb B^{1/p}_{p,q}(B)$ e $n=2^r\ge2$, vale $\Rfun{\Library,H_B}{n}(f)_p\le C_2(B+1)\,n^{-1/p}$.
 \item Para todo $f\in\mathbb B^{1/p}_{p,q}(B)$, vale $\big\|\big(2^{r/p}\Rfun{\Library,H_B}{2^r}(f)_p\big)_{r\ge1}\big\|_{\ell^q}\le C_3(B+1)$.
 \item Se $\|f\|^*_{\Bes{1/p}{p}{q}}\le B$, então $f\in\Rstr{1/p}{q}{\Library,H_B}_p$. Reciprocamente, $\Rstr{1/p}{q}{\Library,H_B}_p\subseteq\Bes{1/p}{p}{q}$ para toda $f\in L^p$. Assim, a bola de raio $B$ de $\Bes{1/p}{p}{q}$ está contida em $\Rstr{1/p}{q}{\Library,H_B}_p$, que está contido em $\Bes{1/p}{p}{q}$.
\end{enumerate}
\end{teorema}

\begin{proof}
Para (a), a desigualdade \eqref{lp:eq:nikolskii} em $m=2$ e a cota de quase-interpolação dão $\sup|Q_2f|\le\mathsf B_0(f)\le C_0B$, com $C_0=8\,2^{1/p}(1+9C_W)$. O \cref{lp:lem:critical-jump} dá saltos $\le J=4C_0B+288C_W2^{1/p}B$. A soma telescópica com $C_p=4$ em \eqref{lp:eq:nikolskii} dá $|\lambda_j|\le M_r$. Tome $C_M=C_0+144C_W2^{1/p}+1$ e $C_J=4C_0+288C_W2^{1/p}+1$. Assim $M_r+1\le C_M(B+1)(r+1)$ e $J+1\le C_J(B+1)$. Como $r+1\le n$, com $A=16C_M+C_J+8$ a agenda satisfaz $8n\le H$, $16(M_r+1)\le H$ e $n(J+1)\le H$, e o \cref{lp:lem:nodal-jump} aplica-se a $Q_nf$. Somando o erro de quase-interpolação (\cref{lp:cor:quasi-error}) e o erro de quantização, obtém-se (a). Para (b), $\mathscr R\le\omega(f,1/N)+5\|f-g\|_p$ (\cref{lp:prop:transfer}), e $\omega(f,1/N)\le BN^{-1/p}$; o termo de quantização é $\le C(r+2)2^{-r}\le C'2^{-r/p}$, pois $(r+2)2^{-r(1-1/p)}$ é limitado. Para (c), a sequência é majorada por uma soma de dois termos: $\omega$-termos em $\ell^q$ com norma $\le B$, e termos de quantização dominados por $C(r+2)\rho^r$ com $\rho=2^{1/p-1}<1$, cuja norma $\ell^q$ é finita. Para (d), a primeira afirmação é o item (c), pois $\Rstr{1/p}{q}{\Library,H_B}_p$ é definido pela norma $\ell^q$ de $(2^{r/p}\Rfun{\Library,H_B}{2^r}(f)_p)_{r\ge1}$. A inclusão reciprocial é o \cref{lp:cor:inclusion}, válido para toda agenda.
\end{proof}

\begin{observacao}
\label{lp:obs:critical-scope}
Para $p>1$ e $0<q<\infty$, o \cref{lp:thm:critical-linear} permite usar $H_B$ em lugar de $\widehat H^{[B]}_A$ na bola de raio $B$ com $s=1/p$. A inclusão direta do item (d) parte dessa bola. A inclusão do estrato em Besov vale em todo $L^p$ para qualquer agenda. Para a igualdade global na diagonal, aplica-se o item (iii) do \cref{lp:conj:besov} com $\widehat H_{q,\kappa}$, que já é linear para $p>1$, $0<q\le1$ e $\kappa=0$. Na diagonal com $p>1$, a extensão da agenda linear a $q=\infty$ e a igualdade global com agenda linear para $1<q<\infty$ ficam fora do alcance deste teorema.
\end{observacao}

\subsection{Piso aritmético e complexidade em bits}
\label{lp:sec:bits}

Escrevemos $e^{H}_{N,p}(f)=\inf\{\|f-g\|_p:g\in\{0\}\cup\Aclass_{N,H}([0,1])\}$. Para $H\ge1$ inteiro, seja $\mathcal Q_H=\{a/b\in\Q:\height(a/b)\le H\}$.

\begin{proposicao}[Piso aritmético uniforme em $L^p$]
\label{lp:prop:floor}
Sejam $H\ge1$ inteiro, $a_H=1/(2H)$ e $f_H\equiv a_H$ em $[0,1]$. Então $\mathrm{dist}(a_H,\mathcal Q_H)=1/(2H)$ e, para todo $N\ge1$ e todo $g\in\{0\}\cup\Aclass_{N,H}([0,1])$,
\begin{equation}
 \|f_H-g\|_p\ge\frac1{2H}.
 \label{lp:eq:floor}
\end{equation}
Em particular, $e^{H}_{N,p}(f_H)\ge1/(2H)$ para todo $N$.
\end{proposicao}

\begin{proof}
Todo racional positivo de $\mathcal Q_H$ tem numerador inteiro positivo e denominador no máximo $H$, logo é pelo menos $1/H$, e dista de $a_H$ ao menos $1/(2H)$. Os racionais não positivos distam de $a_H$ ao menos $a_H=1/(2H)$, e $0\in\mathcal Q_H$ realiza essa distância.

Para $g=0$, $\|f_H\|_p=a_H$. Seja $g\in\Aclass_{N,H}([0,1])$ e seja $J$ uma folha da representação, com $g(x)=y_0+s(x-x_J)$ em $J$, onde $x_J$ é o centro de $J$ e $\height(y_0)\le H$. Como $g$ é afim, $\int_Jg=|J|\,y_0$. A desigualdade de Jensen para a função convexa $t\mapsto|t|^p$ dá
\[
 \int_J|f_H-g|^p\ge|J|\Bigl|\frac1{|J|}\int_J(f_H-g)\Bigr|^p=|J|\,|a_H-y_0|^p\ge|J|\,(2H)^{-p},
\]
pois $y_0\in\mathcal Q_H$. As folhas têm interiores disjuntos e união $[0,1]$. Somando, $\|f_H-g\|_p^p\ge(2H)^{-p}$.
\end{proof}

\begin{corolario}[Necessidade do expoente aritmético]
\label{lp:cor:floor}
Sejam $s>0$, $C>0$ e uma agenda $H$ com $H(N)\ge2$. Suponha que, para algum $N\ge2$,
\[
 \sup_{f\in\mathbb B^s_{p,q}(1)}\Rfun{\Library,H}{N}(f)_p\le C\,N^{-s}.
\]
Então $H(N)\ge N^s/(2C)$.
\end{corolario}

\begin{proof}
Como $H(N)\ge2$, a constante $f\equiv a_{H(N)}$ tem $a_{H(N)}\le1/4$, e portanto $\|f\|^*_{\Bes{s}{p}{q}}=a_{H(N)}\le1$, pois seu módulo de segunda ordem é nulo. Para todo $g\in\Sig{N}$, $\|f-g\|_p+\wt{g}{N^{-1}}{p}\ge\|f-g\|_p\ge1/(2H(N))$ pela \cref{lp:prop:floor}. Logo $1/(2H(N))\le CN^{-s}$.
\end{proof}

\begin{observacao}
\label{lp:obs:floor}
Sejam $s\ge1$ com $s>1/p$ e $\kappa=0$. Pela prova do \cref{lp:thm:uniform}, para $n=2^r$ vale
\[
 \sup_{f\in\mathbb B^s_{p,q}(B)}\Rfun{\Library,\widehat H}{n}(f)_p
 \le\bigl((1+45C_W)B+\tfrac52\bigr)n^{-s},
\]
com $\widehat H(N)\asymp(B+1)N^s$. O \cref{lp:cor:floor} dá a potência necessária nessa faixa para a taxa uniforme por escala. Para $s<1$, o piso geométrico é mais forte que o aritmético, e para $s<1/p$ o piso por inclinação melhora ainda essa cota. O \cref{lp:cor:joint-necessity} aplica os três pisos à biblioteca inteira, incluindo as partições adaptativas.
\end{observacao}

Para a codificação binária, o orçamento de altura $H$ é conhecido pelo codificador e pelo decodificador. Um bit indica $g=0$ ou $g\ne0$. Para $g\ne0$, escolhe-se uma representação admissível. A forma da árvore, em pré-ordem, usa um bit por nó para distinguir folha de nó interno e determina o par $(k,m)$ de cada folha. Cada folha traz as duas coordenadas racionais $(y_0,s)$, com sinal, numerador e denominador.

\begin{proposicao}[Bits de $\Aclass_{N,H}$]
\label{lp:prop:bits}
Sejam $N\ge1$ e $H\ge2$ inteiros e $c_H=\lceil\log_2(H+1)\rceil$. Toda $g\in\{0\}\cup\Aclass_{N,H}([0,1])$ admite uma descrição binária livre de prefixo com no máximo $4N(c_H+1)$ bits. Em particular, $\log_2\#\bigl(\{0\}\cup\Aclass_{N,H}([0,1])\bigr)\le4N(c_H+1)$.
\end{proposicao}

\begin{proof}
Uma coordenada $a/b$ reduzida, com $b>0$ e $\height(a/b)\le H$, tem $|a|\le H<2^{c_H}$ e $1\le b\le H<2^{c_H}$. Ela usa então $2c_H+1$ bits, com um de sinal e $c_H$ para cada inteiro, em comprimento fixo. Uma folha usa $4c_H+2$ bits. Uma árvore com $N'\le N$ folhas tem $2N'-1$ nós. O comprimento total é, portanto, no máximo $1+(2N-1)+N(4c_H+2)=4N(c_H+1)$. A forma da árvore em pré-ordem é livre de prefixo, e o número de folhas, que ela determina, fixa o comprimento do restante. Escolhida uma descrição por função, completar cada palavra com zeros até $L=4N(c_H+1)$ preserva a injetividade, pois duas palavras com o mesmo preenchimento seriam comparáveis por prefixo. Há $2^L$ palavras desse comprimento.
\end{proof}

\begin{corolario}[Taxa e bits sobre bolas]
\label{lp:cor:bits}
Sejam $1\le p<\infty$, $0<s<2$, $0<q\le\infty$ e $\widehat H=\widehat H^{[B]}_{A}$ com $\kappa=0$ e $A$ como no \cref{lp:thm:uniform}. Existe $c=c(p,s,q)$ tal que, para todo $B\ge1$, todo $f\in\mathbb B^s_{p,q}(B)$ e todo $n=2^r\ge2$, há $g\in\Sig{n}$ com
\[
 \|f-g\|_p\le(9C_WB+1)\,n^{-s}
\]
e descrição binária de comprimento no máximo $c\,n\,\bigl(\log n+\log(B+1)\bigr)$.
\end{corolario}

\begin{proof}
O \cref{lp:thm:uniform}, com $\kappa=0$, fornece $g\in\Aclass_{n,\widehat H(n)}([0,1])$ satisfazendo a cota de erro. Para $\kappa=0$, vale $\widehat H(n)\le A(B+1)(n+1)^{\hat\nu}\mathsf L(n)^{\theta_q}+1$, com $\hat\nu\le2$ e $\theta_q\le1$. Logo $c_{\widehat H(n)}=O(\log n+\log(B+1))$, com constante dependente apenas de $(p,s,q)$, e a \cref{lp:prop:bits} conclui.
\end{proof}

\begin{observacao}
\label{lp:obs:entropy}
Seja $\mathcal N(\varepsilon)$ o número mínimo de bolas de raio $\varepsilon$ em $L^p([0,1])$ que cobrem $\mathbb B^s_{p,q}(B)$. O teorema de Birman--Solomyak \citep[Teorema~1]{lp-kerkyacharianpicard2003}, com dimensão $1$ e o mesmo expoente $p$ no espaço de origem e na norma do erro, dá $\log_2\mathcal N(\varepsilon)\asymp (B/\varepsilon)^{1/s}$ para $\varepsilon/B$ suficientemente pequeno, por equivalência das quase-normas de Besov e homogeneidade. Vale para $1\le p<\infty$, $0<s<2$ e $0<q\le\infty$.

Com $B$ fixo e erro uniforme $\varepsilon=C_Bn^{-s}$, as saídas de qualquer decodificador com comprimento máximo $L$ formam uma cobertura por menos de $2^{L+1}$ centros, logo $L\gtrsim n$. Se as saídas podem ser centros arbitrários em $L^p$, indexar uma cobertura de raio $\varepsilon$ também dá $O(n)$ bits. O formato explícito do \cref{lp:cor:bits} usa $O(n\log n)$ bits para reconstruir os aproximantes da biblioteca. Portanto seu custo fica a um fator de, no máximo, ordem $\log n$ do mínimo sem restrição sobre os centros. Esse comprimento mede a descrição de um aproximante já escolhido. A questão de \emph{calcular} um aproximante ótimo por programação dinâmica é tratada na Parte~III de \citet{lp-oliveiravol} (Teoremas III.9.1 e III.9.3) para o erro em norma uniforme.
\end{observacao}

\subsection{Pisos geométrico e de inclinação, necessidade e transferência}
\label{lp:sec:floors}

Os pisos a seguir completam a \cref{lp:prop:floor}. Eles valem para a biblioteca completa $\Sig N$, não apenas para o piso aritmético, e as constantes dependem só de $(p,s,q)$.

\begin{lema}[Escala de um bump]
\label{lp:lem:bump}
Sejam $1\le p<\infty$, $0<s<2$ e $0<q\le\infty$. Fixe $\varphi\in C^\infty_c((1/8,5/8))$ com $\varphi=1$ em $[1/4,1/2]$, e ponha $M_\varphi=\sup|\varphi''|$ e $\kappa_\varphi=\int_0^1(\varphi')^2>0$. Existe $\Lambda\ge1$, dependendo só de $(p,s,q,\varphi)$, tal que, para $a\in\mathbb R$ e $0<t\le1$, a função $\psi_{a,t}(x)=a\,\varphi\bigl((x-1+t)/t\bigr)$ satisfaz
\[
 \|\psi_{a,t}\|^*_{\Bes sp q}\le\Lambda\,|a|\,t^{1/p-s}.
\]
Além disso, para $a\ge0$, $0<t\le1$ e toda função afim $\ell$,
\[
 \|\psi_{a,t}-\ell\|_{L^p(1-t,1)}\ge\frac{\kappa_\varphi}{M_\varphi}\,a\,t^{1/p}.
\]
\end{lema}

\begin{proof}
Para a norma, $\|\psi_{a,t}\|_p=|a|t^{1/p}\|\varphi\|_p$, e $\psi_{a,t}''=a\,t^{-2}\varphi''(\cdot)$. A fórmula $\Delta_h^2\psi(x)=\int_0^h\!\int_0^h\psi''(x+u+v)\,du\,dv$, a desigualdade de Minkowski e a subaditividade dão
\[
 \wt{\psi_{a,t}}{h}{p}\le|a|t^{1/p}\,\Omega(h/t),\qquad \Omega(u)=\min\{4\|\varphi\|_p,\;u^2\|\varphi''\|_p\}.
\]
Para $h_r=2^{-r}$ e $u_r=h_r/t$, temos $|a|t^{1/p}2^{rs}=|a|t^{1/p-s}\,u_r^{-s}$, logo
\[
 2^{rs}\wt{\psi_{a,t}}{h_r}{p}\le|a|t^{1/p-s}\,u_r^{-s}\Omega(u_r).
\]
Como $u_{r+1}=u_r/2$, para $q<\infty$ obtemos as cotas geométricas
\[
 \sum_{u_r\ge1}u_r^{-sq}\le(1-2^{-sq})^{-1},
 \qquad \sum_{u_r<1}u_r^{(2-s)q}\le(1-2^{-(2-s)q})^{-1}.
\]
Ambas são finitas porque $0<s<2$. Para $q=\infty$, usa-se que $u^{-s}\Omega(u)$ é limitada. Assim a sequência tem norma $\ell^q$ limitada por $|a|t^{1/p-s}K_1$, com $K_1$ dependendo só de $(p,s,q,\varphi)$. Como $\|\psi_{a,t}\|_p\le|a|t^{1/p-s}\|\varphi\|_p$, basta tomar $\Lambda=\max\{1,\|\varphi\|_p+K_1\}$.

Para a segunda desigualdade, seja $L(y)=\ell(1-t+ty)$, afim. Integrando por partes, com $\varphi(0)=\varphi(1)=\varphi'(0)=\varphi'(1)=0$, obtemos $\int_0^1\varphi\,\varphi''=-\kappa_\varphi$ e $\int_0^1 L\,\varphi''=0$. Logo $a\kappa_\varphi\le M_\varphi\int_0^1|a\varphi(y)-L(y)|\,dy$. Pela substituição $x=1-t+ty$, $\int_{1-t}^1|\psi_{a,t}-\ell|\,dx=t\int_0^1|a\varphi-L|\,dy\ge t\,a\kappa_\varphi/M_\varphi$. Como o intervalo $(1-t,1)$ tem comprimento $t$, a desigualdade de Hölder dá $\int_{1-t}^1|h|\le t^{1-1/p}\|h\|_{L^p(1-t,1)}$, e portanto $\|\psi_{a,t}-\ell\|_{L^p(1-t,1)}\ge t^{1/p-1}\cdot t\,a\kappa_\varphi/M_\varphi=\frac{\kappa_\varphi}{M_\varphi}a\,t^{1/p}$.
\end{proof}

\begin{proposicao}[Piso geométrico]
\label{lp:prop:geom-floor}
Sejam $1\le p<\infty$, $0<s<2$ e $0<q\le\infty$. Existe $c_G=c_G(p,s,q)>0$ tal que, para todo inteiro $H\ge2$, existe $f_H$ mensurável com $\|f_H\|^*_{\Bes sp q}\le1$ e
\[
 \|f_H-g\|_p\ge c_G\,H^{-s}\qquad\text{para todo }N\ge1\text{ e todo }g\in\Sig N\text{ com agenda }H.
\]
\end{proposicao}

\begin{proof}
Seja $\Lambda$ a constante do \cref{lp:lem:bump}, com $\varphi$ fixo. Tome $t=1/H$ e $a=t^{s-1/p}/\Lambda$, e seja $f_H=\psi_{a,t}$. Pelo \cref{lp:lem:bump}, $\|f_H\|^*\le\Lambda a t^{1/p-s}=1$. Seja $g\in\Sig N$ com agenda $H$. Pela última folha, $g$ é afim em $[1-1/H,1]=[1-t,1]$, e como $f_H=0$ fora de $(1-t,1)$, o mesmo lema dá
\[
 \|f_H-g\|_p\ge\|f_H-g\|_{L^p(1-t,1)}\ge\frac{\kappa_\varphi}{M_\varphi}\,a\,t^{1/p}=\frac{\kappa_\varphi}{M_\varphi\Lambda}\,t^{s}=c_G\,H^{-s},
\]
com $c_G=\kappa_\varphi/(M_\varphi\Lambda)$, que depende só de $(p,s,q)$.
\end{proof}

\begin{proposicao}[Piso por inclinação]
\label{lp:prop:slope-floor}
Sejam $1\le p<\infty$, $0<s\le1/p$ e $0<q\le\infty$, e ponha $\beta=s/(1+1/p-s)$. Existe $c_L=c_L(p,s,q)>0$ tal que, para todo inteiro $H\ge2$, existe $f_H$ mensurável com $\|f_H\|^*_{\Bes sp q}\le1$ e
\[
 \|f_H-g\|_p\ge c_L\,H^{-\beta}\qquad\text{para todo }N\ge1\text{ e todo }g\in\Sig N\text{ com agenda }H.
\]
\end{proposicao}

\begin{proof}
Ponha $\delta=1/p-s$ e $c=\min\{\Lambda^{-1},2^{-(1+\delta)}\}$, com $\Lambda$ do \cref{lp:lem:bump}. Tome $t=(c/(2H))^{1/(1+\delta)}\le1/2$, $a=c\,t^{-\delta}$ e $f_H(x)=a\varphi((x-1/4)/t)$. O suporte fica em $(0,1)$. As cotas de $\|f_H\|_p$ e do módulo no \cref{lp:lem:bump} valem com a mesma constante para esse suporte transladado, pois usam as normas de $\varphi$ e $\varphi''$. Assim $\|f_H\|^*\le\Lambda a t^\delta\le1$.

As funções $g$ são $H$-Lipschitz em $[0,1]$. Para $x$ em $I=[1/4+t/4,1/4+t/2]$, temos $f_H(x)=a$ e $f_H(x+t)=0$. Como $Ht=a/2$, segue $|f_H-g|(x)+|f_H-g|(x+t)\ge a/2$. Os intervalos $I$ e $I+t$ são disjuntos e estão em $[0,1]$. Por convexidade, a soma das potências $p$ dos dois erros é ao menos $2^{1-p}(a/2)^p$. Integrando em $I$, de comprimento $t/4$, obtemos
\[
 \|f_H-g\|_p\ge 2^{-2-1/p}a\,t^{1/p}
 =2^{-2-1/p}c(c/2)^\beta H^{-\beta}.
\]
\end{proof}

\begin{corolario}[Necessidade conjunta]
\label{lp:cor:joint-necessity}
Sejam $1\le p<\infty$, $0<s<2$, $0<q\le\infty$, $C>0$, $N\ge1$ e uma agenda $H$ com $H(N)\ge2$. Suponha que, para todo $f$ mensurável com $\|f\|^*_{\Bes sp q}\le1$,
\[
 e^{H(N)}_{N,p}(f)\le C\,N^{-s}.
\]
Então $H(N)\ge(c_G/C)^{1/s}N$. Se $s\le1/p$, então $H(N)\ge(c_L/C)^{1/\beta}N^{1+1/p-s}$. Além disso, $H(N)\ge N^s/(2C)$. Se a hipótese for feita com $\mathscr R$ no lugar de $e$, ela implica a hipótese acima, porque $e\le\mathscr R$ pela \cref{lp:prop:transfer}; logo as mesmas conclusões valem.
\end{corolario}

\begin{proof}
Seja $H=H(N)$. Pela \cref{lp:prop:geom-floor}, existe $f_H$ com $\|f_H\|^*\le1$ tal que $\|f_H-g\|_p\ge c_GH^{-s}$ para todo $g\in\{0\}\cup\Aclass_{N,H}([0,1])$, o conjunto do ínfimo que define $e^H_{N,p}$. Como $0$ é afim e cada $g\in\Aclass_{N,H}([0,1])$ é afim em $[1-1/H,1]$ pela última folha, a proposição se aplica a esse conjunto. Logo $c_GH^{-s}\le e^H_{N,p}(f_H)\le CN^{-s}$, o que dá $H^s\ge(c_G/C)N^s$. Se $s\le1/p$, a \cref{lp:prop:slope-floor} dá $c_LH^{-\beta}\le CN^{-s}$, isto é, $H\ge(c_L/C)^{1/\beta}N^{s/\beta}$, com $s/\beta=1+1/p-s$. Para a última desigualdade, $f\equiv1/(2H)$ tem $\|f\|^*=1/(2H)\le1$, e pela \cref{lp:prop:floor} $e^H_{N,p}(f)\ge1/(2H)$; logo $1/(2H)\le CN^{-s}$.
\end{proof}

\begin{proposicao}[Transferência exata entre erro e realização]
\label{lp:prop:transfer}
Sejam $1\le p<\infty$, $G_N$ famílias de funções mensuráveis e $f$ mensurável. Com $e_N(f)=\inf_{g\in G_N}\|f-g\|_p$ e $\mathscr R_N(f)=\inf_{g\in G_N}\{\|f-g\|_p+\omega_p(g,1/N)\}$,
\[
 e_N(f)\le\mathscr R_N(f)\le\omega_p(f,1/N)+5\,e_N(f).
\]
\end{proposicao}

\begin{proof}
A primeira desigualdade vale porque $\mathscr R_N$ soma $\omega_p(g,1/N)\ge0$ a cada termo do ínfimo. Para a segunda, se $g\in G_N$, então $\omega_p(g,1/N)\le\omega_p(f,1/N)+4\|f-g\|_p$ pela subaditividade de $\omega_p$ e pela cota $\omega_p(h,t)\le4\|h\|_p$. Logo $\mathscr R_N(f)\le\omega_p(f,1/N)+5\|f-g\|_p$ para todo $g$, e o ínfimo sobre $g$ dá a cota.
\end{proof}

\begin{observacao}
\label{lp:obs:floor-closure}
A \cref{lp:cor:joint-necessity} determina a potência necessária da altura para a taxa uniforme por escala $CN^{-s}$ sobre a bola unitária, tanto para o melhor erro em $L^p$ quanto para a realização. O piso por inclinação dá $H(N)\gtrsim N^{1+1/p-s}$ se $s<1/p$, o geométrico dá $H(N)\gtrsim N$ se $1/p\le s\le1$ e o aritmético dá $H(N)\gtrsim N^s$ se $1<s<2$. As constantes dessas cotas dependem de $(p,s,q,C)$. O alvo de teste pode variar com $H(N)$, e uma cota uniforme para uma bola de raio $B\ge1$ implica a mesma necessidade porque a bola contém a unitária.

As cotas superiores do \cref{lp:thm:uniform}, com $\kappa=0$, e do \cref{lp:thm:critical-linear} atingem essas potências nos orçamentos diádicos. Para $B$ fixo, quando a cota superior da \cref{lp:tab:schedules} inclui $\mathsf L_1^{\theta_q}$, a comparação determina apenas a potência de $N$. A necessidade desse fator na taxa por escala e a necessidade de fatores adicionais para somabilidade são questões distintas. Os pisos também deixam indeterminada a altura exigida para um alvo fixo.
\end{observacao}

\subsection{Agendas de log iterado}
\label{lp:sec:loglog}

Ponha $\mathsf L_1(N)=\log(eN)$ e $\mathsf L_2(N)=\log\bigl(e\,\mathsf L_1(N)\bigr)=1+\log\mathsf L_1(N)$. Ambas são não decrescentes e valem pelo menos $1$ para $N\ge1$. Para $0<q<\infty$, $\gamma>0$ e $s\ge1$, a \emph{agenda de log iterado} é
\begin{equation}
 H^{(2)}_{s;q,\gamma}(N)=\Bigl\lceil8\,(N+1)^{s}\,\mathsf L_1(N)^{1/q}\,\mathsf L_2(N)^{\gamma}\Bigr\rceil .
 \label{lp:eq:loglog}
\end{equation}

\begin{proposicao}[Critério exato para agendas regulares]
\label{lp:prop:regular}
Sejam $s\ge1$, $0<q<\infty$ e $h\colon\N\to[1,\infty)$ não decrescente, e ponha $H(N)=\lceil8(N+1)^{s}h(N)\rceil$. Então
\[
 \Bigl(\frac{2^{rs}}{H(2^r)}\Bigr)_{r\ge1}\in\ell^q
 \iff\sum_{r\ge1}h(2^r)^{-q}<\infty
 \iff\sum_{N\ge1}\frac1{N\,h(N)^q}<\infty .
\]
\end{proposicao}

\begin{proof}
Como $h\ge1$, vale $8\,2^{rs}h(2^r)\le H(2^r)\le(8\cdot2^s+1)\,2^{rs}h(2^r)$, de modo que $2^{rs}/H(2^r)$ é comparável a $1/h(2^r)$, com constantes que dependem só de $s$. Isso prova a primeira equivalência. A sequência $t_N=1/(Nh(N)^q)$ é não crescente, e $2^rt_{2^r}=h(2^r)^{-q}$. O teste de condensação de Cauchy dá a segunda equivalência.
\end{proof}

\begin{corolario}[Agenda de log-log]
\label{lp:cor:loglog}
Sejam $1\le p<\infty$, $1\le s<2$ com $s>1/p$, $0<q<\infty$ e $\gamma>1/q$. Então
\[
 \Bes{s}{p}{q}=\Rstr{s}{q}{\Library,H^{(2)}_{s;q,\gamma}}_p .
\]
O mesmo vale, com a mesma agenda, no \cref{lp:conj:stability} e no \cref{lp:thm:approx-stability}. Além disso, o \cref{lp:thm:uniform} vale com a agenda uniforme $\lceil A(B+1)(N+1)^s\,\mathsf L_1(N)^{1/q}\mathsf L_2(N)^{\gamma}\rceil$ no lugar de \eqref{lp:eq:uniform-schedule}, e com $c_2$ dependendo de $(p,s,q,\gamma)$.
\end{corolario}

\begin{proof}
A função $h=\mathsf L_1^{1/q}\mathsf L_2^{\gamma}$ é não decrescente e $\ge1$, e $H^{(2)}_{s;q,\gamma}(N)\ge8N^{\nu_s}$. Resta verificar o critério \eqref{lp:eq:summability}, isto é, pela \cref{lp:prop:regular}, que $\sum_rh(2^r)^{-q}<\infty$. Para $r\ge2$, $\mathsf L_1(2^r)=1+r\log2\ge r\log2$ e $\mathsf L_2(2^r)=1+\log(1+r\log2)\ge\log r$, pois $1+\log\log2>0$. Logo
\[
 h(2^r)^{-q}=\mathsf L_1(2^r)^{-1}\mathsf L_2(2^r)^{-\gamma q}\le\frac1{\log2}\cdot\frac1{r(\log r)^{\gamma q}},
\]
e $\sum_{r\ge3}r^{-1}(\log r)^{-a}$ converge para $a=\gamma q>1$, pelo teste da integral. A \cref{lp:thm:besov-conditional} dá a igualdade. As afirmações sobre estabilidade seguem da \cref{lp:thm:stability-conditional} e da \cref{lp:thm:approx-transfer}, que dependem da agenda apenas por \eqref{lp:eq:summability}. Para a versão uniforme, basta trocar $\mathsf L(N)^\kappa$ por $\mathsf L_1(N)^{1/q}\mathsf L_2(N)^{\gamma}$ em $Q(n)$ e em \eqref{lp:eq:uniform-schedule}. As contas do \cref{lp:lem:uniform-height} permanecem, e a sequência $2^{rs}/Q(2^r)$ é limitada por $h(2^r)^{-1}$, que pertence a $\ell^q$ pelo argumento acima.
\end{proof}

\begin{observacao}
\label{lp:obs:loglog}
Para todo $\kappa>1/q$, vale $H^{(2)}_{s;q,\gamma}(N)/H_{s,\kappa}(N)\to0$, pois $\mathsf L_2^{\gamma}/\mathsf L_1^{\kappa-1/q}\to0$. A agenda de log-log é, portanto, estritamente menor que qualquer agenda com expoente fixo $\kappa>1/q$ do \cref{lp:conj:besov}. O mesmo argumento com $m$ logaritmos, $\mathsf L_1^{1/q}\cdots\mathsf L_{m-1}^{1/q}\mathsf L_m^{\gamma}$ e $\gamma>1/q$, dá a série de Bertrand correspondente. Pela \cref{lp:prop:regular}, a somabilidade do peso $2^{rs}/H(2^r)$ equivale à convergência de $\sum_rh(2^r)^{-q}$. Para $\gamma=1/q$, a série diverge por comparação com $\sum_r1/(r\log r)$. O limiar se refere ao peso usado na prova da inclusão direta. Os pisos do \cref{lp:cor:joint-necessity} dão a potência necessária à taxa uniforme por escala, e a necessidade dos fatores logarítmicos para a caracterização permanece indeterminada.

O mesmo vale no item (iii) do \cref{lp:conj:besov} com $p=1$, $s=1$ e $q<\infty$: se $Q(n)$ é a menor potência de dois com $Q(n)\ge n\,\mathsf L_1(n)^{1/q}\mathsf L_2(n)^{\gamma}$ e $\gamma>1/q$, a quantização contribui com $n^{s}/Q(n)\le(\mathsf L_1^{1/q}\mathsf L_2^{\gamma})^{-1}$, e a agenda passa a ser $\lceil8(N+1)\,\mathsf L_1^{\theta_q+1/q}\mathsf L_2^{\gamma}\rceil$.
\end{observacao}

A \cref{lp:tab:schedules} reúne as agendas suficientes para a taxa uniforme nos orçamentos diádicos $N=2^r\ge2$ e para a somabilidade do funcional de realização.

\begin{table}[H]
\centering
\caption{Alturas suficientes para a taxa por escala e para a realização}
\label{lp:tab:schedules}
\small
\tablecolors
\begin{tabularx}{\linewidth}{L{0.25\linewidth}L{0.15\linewidth}L{0.20\linewidth}>{\raggedright\arraybackslash}X}
\toprule
\tablehead{Regularidade} & \tablehead{Índice $q$} & \tablehead{Taxa $N^{-s}$ na bola} & \tablehead{Realização em $\ell^q$} \\
\midrule
$0<s<1/p$ & $0<q\le\infty$ & $N^{1+1/p-s}$ & $N^{1+1/p-s}$ \\
$s=1/p$, $p>1$ & $0<q<\infty$ & $N$ & $N$ na bola ou $N\mathsf L_1^{\theta_q}$ global \\
$s=1/p$ & $q=\infty$ & $N\mathsf L_1$ & $N\mathsf L_1$ \\
$s=p=1$ & $0<q<\infty$ & $N\mathsf L_1^{\theta_q}$ & $N\mathsf L_1^{\theta_q+\kappa}$, $\kappa>1/q$ \\
$1/p<s<1$ & $0<q\le\infty$ & $N$ & $N$ \\
$1\le s<2$, $s>1/p$ & $q=\infty$ & $N^s$ & $N^s$ \\
$1\le s<2$, $s>1/p$ & $0<q<\infty$ & $N^s$ & $N^s\mathsf L_1^\kappa$, $\kappa>1/q$ \\
\bottomrule
\end{tabularx}
\par\raggedright
\tablecredit
\par\smallskip\noindent
A taxa na bola usa $B\ge1$ e o fator $B+1$ na altura, incluindo a entrada linear na última coluna. As demais entradas dessa coluna são agendas de caracterização global. Os prefatores são os dos teoremas correspondentes. Na diagonal, $\theta_q$ vem do controle dos valores nodais. Para $s\ge1$ e $q<\infty$, pode-se substituir $\mathsf L_1^\kappa$ por $\mathsf L_1^{1/q}\mathsf L_2^\gamma$ com $\gamma>1/q$. Para $s>1/p$ e $q<\infty$, $\widetilde H_{s,\eta}$ com $\eta>0$ também dá a caracterização global.
\end{table}

\section{Considerações finais}
A somabilidade diádica do funcional de realização caracteriza os espaços de Besov estudados sob os orçamentos simultâneos de folhas e altura aritmética. A estimativa de Jackson fornece a inclusão direta em cada regime de regularidade, e a desigualdade inversa decorre das propriedades do módulo de suavidade. A caracterização se transfere a famílias aninhadas de redes ReLU sob inclusão exata ou aproximada, com o controle de tolerância enunciado no \cref{lp:thm:approx-stability}.

Para a taxa uniforme $N^{-s}$ por escala, os pisos e as cotas superiores fixam a potência da altura em toda a faixa $0<s<2$. A potência é $1+1/p-s$ abaixo da diagonal, $1$ de $s=1/p$ até $s=1$ e $s$ acima de $1$. Na diagonal com $p>1$ e $0<q<\infty$, a agenda linear $H_B$ também dá a cota uniforme em $\ell^q$ e a caracterização na bola de raio $B$. Nos demais casos em que as agendas suficientes incluem logaritmos, os pisos deixam em aberto se esses fatores são necessários. As agendas de log iterado refinam as condições suficientes de somabilidade.

Para raio fixo, a codificação explícita dos aproximantes tem comprimento máximo $O(N\log N)$ na taxa uniforme em $L^p$. A entropia da bola dá uma cota inferior de ordem $N$ para qualquer codificação com o mesmo erro uniforme.
\ifbesovabnt
\articlereferences
\fi

\articleappendix
\section{Formalização das demonstrações em Lean 4}
\label{lp:sec:lean}

O projeto \code{lean4/formalizacao} verifica, com Lean 4.28.0 e Mathlib v4.28.0 \citep{lp-mathlib,lp-lean4}, o \cref{lp:conj:besov}, nos itens (i), (ii) e (iii), o \cref{lp:conj:stability} e o \cref{lp:thm:approx-stability}. A desigualdade inversa é demonstrada para módulos a valores em $[0,\infty]$ e especializada a $L^p$. A prova direta combina Whitney, quase-interpolação, imersão contínua (regime $s>1/p$), cota em sup dos quase-interpolantes (regime $s\le1/p$) e construção da árvore diádica quantizada. A estimativa de Jackson resultante é aplicada às agendas dos itens (i), (ii) e (iii) do \cref{lp:conj:besov} e às escalas deslocadas dos \cref{lp:conj:stability,lp:thm:approx-stability}. A formalização inclui as cotas uniformes sobre bolas, o piso aritmético, as descrições binárias e as agendas de log iterado da \cref{lp:sec:uniform}.

\subsection{Representação dos objetos e dos teoremas}
\label{lp:sec:fidelity}

O código representa as classes de $L^p$ por funções mensuráveis $\R\to\R$ e usa a medida de Lebesgue restrita a $[0,1]$. Em \code{wtP}, a medida é restrita a $[0,1-2h]$, conforme \eqref{lp:eq:modulus-p}. As definições \code{besovP}, \code{realP} e \code{stratumP} correspondem à condição de Besov, ao funcional e ao espaço de realização. A sequência de realização usa o orçamento $2^{r+1}$ para $r\ge0$.

A forma de Jackson com $C_0=9C_W(p)$ e $C_1=16$ é \code{hypJpStrong_of_lt}, aplicada ao representante contínuo fornecido por \code{besov_continuous_rep}. As reduções usam \code{HypJp}, demonstrada em \code{hypJp_of_lt}, com constantes que podem depender de $f$. Os teoremas finais do regime $s>1/p$ incorporam essa prova e usam $1/p<s$. No artigo, $s<2$ é o intervalo da caracterização clássica de Besov por diferenças de ordem dois.

Para $s\le1/p$, \code{jackson_low} (\cref{lp:thm:jackson-low}) tem constantes $C_W$ e $K$ que dependem de $f$, em vez de $C_0(S+1)$, e \code{besov_subset_stratum_low} (\cref{lp:thm:low-transfer}) usa o deslocamento diádico de $f$ por meio da condição \code{hshift}. Os enunciados \code{conj1_low}, \code{conj1_diag}, \code{conj2_low} e \code{conj2_diag} escrevem $p\in[1,\infty)$ como número real, e as agendas são \texttt{Hsk (1 + 1/p - s) 0} e \texttt{Hsk 1 (thetaQ q + kappa)}, que coincidem com \eqref{lp:eq:schedule-low} e \eqref{lp:eq:schedule-diag}. O item (ii) em Lean não exige $s>0$.

A estabilidade no código se aplica a famílias de funções mensuráveis sob \eqref{lp:eq:nesting} e \eqref{lp:eq:contains} (ou \eqref{lp:eq:approx-contains} com \eqref{lp:eq:approx-summable}). Como as famílias consideradas no texto são subconjuntos de $C([0,1])$, que são mensuráveis, essa exigência não restringe os enunciados dos \cref{lp:conj:stability,lp:thm:approx-stability}. Na estabilidade aproximada, a tolerância é uma sequência $\varepsilon\colon\N\to[0,\infty]$ (\texttt{eps : N -> ENNReal}), e a hipótese de somabilidade \eqref{lp:eq:approx-summable} entra como premissa de \code{Conj2ApproxStability}. O teorema \code{besov_subset_stratum_fam} do regime $s>1/p$ é o caso $\varepsilon=0$ de seu núcleo com tolerância, e o mesmo vale para \code{besov_subset_stratum_low} no regime $s\le1/p$. O módulo \code{Stability} transfere a estimativa de Jackson para o orçamento diádico deslocado.

\Cref{lp:tab:verified} localiza as declarações para as demonstrações principais.

\begin{leantable}{Demonstrações e declarações correspondentes em Lean}{lp:tab:verified}{Resultado}{Declarações Lean}
Desigualdade inversa e inclusão em Besov & \leanname{wtP_le_four_realP}, \leanname{stratum_subset_besov_conn} \\
Projeção afim e quase-interpolação & \leanname{kernel_is_affine_projection}, \leanname{local_error_inf}, \leanname{lam_diff_le} \\
Whitney e erro global & \leanname{whitney_lp}, \leanname{quasi_error_unc} \\
Imersão e módulo de primeira ordem & \leanname{besov_continuous_rep}, \leanname{modulus_lt1_unc}, \leanname{modulus_eq1_unc}, \leanname{modulus_gt1_unc} \\
Quantização nodal e Jackson ($s>1/p$) & \leanname{nodal_tree}, \leanname{lp_jackson_core}, \leanname{hypJpStrong_of_lt} \\
Caracterização de Besov ($s>1/p$) & \leanname{conj1_endpoint}, \leanname{conj1_subcritical}, \leanname{conj1_critical}, \leanname{conj1_tilde} \\
Estabilidade por inclusão ($s>1/p$) & \leanname{conj2_endpoint}, \leanname{conj2_subcritical}, \leanname{conj2_critical}, \leanname{conj2_tilde} \\
Quantização sem controle dos saltos, cota em sup e Jackson ($s\le1/p$) & \leanname{nodal_tree_gen}, \leanname{lam_bound}, \leanname{jackson_low} \\
Transferência e caracterização ($s\le1/p$) & \leanname{besov_subset_stratum_low}, \leanname{conj1_low}, \leanname{conj1_diag} \\
Estabilidade por inclusão ($s\le1/p$) & \leanname{conj2_low}, \leanname{conj2_diag} \\
Estabilidade aproximada ($s>1/p$) & \leanname{conj2_approx_endpoint}, \leanname{conj2_approx_subcritical}, \leanname{conj2_approx_critical}, \leanname{conj2_approx_tilde} \\
Estabilidade aproximada ($s\le1/p$) & \leanname{besov_subset_stratum_low_approx}, \leanname{conj2_approx_low}, \leanname{conj2_approx_diag} \\
\end{leantable}

\par\ifbesovabnt\Needspace{8\baselineskip}\else\Needspace{14\baselineskip}\fi
\Cref{lp:tab:uniform-lean} localiza as declarações para as extensões da \cref{lp:sec:uniform}.

\begin{leantable}{Declarações Lean para as extensões da Seção~6}{lp:tab:uniform-lean}{Resultado}{Declarações Lean}
Quantização e cotas uniformes & \leanname{uniform_height_quantization}, \leanname{lam_bound_unif}, \leanname{uniform_quantWeight_lt_top}, \leanname{uniform_besov_ball} \\
Piso aritmético em $L^p$ & \leanname{arithmetic_floor_lp}, \leanname{arithmetic_floor_bestError}, \leanname{arithmetic_height_necessary} \\
Código binário e taxa & \leanname{boundedLibrary_code}, \leanname{boundedLibrary_log_card}, \leanname{uniform_rate_bits} \\
Critério exato da agenda & \leanname{regular_schedule_criterion}, \leanname{iterated_schedule_iff} \\
Log-log e estabilidade & \leanname{conj1_loglog}, \leanname{conj2_loglog}, \leanname{conj2_approx_loglog}, \leanname{uniform_loglog_ball}, \leanname{conj1_loglog_diagonal} \\
\end{leantable}

\Needspace{9\baselineskip}
\subsection{Pisos, necessidade e altura crítica}
\label{lp:sec:lean-floors}

\Cref{lp:tab:critical-lean} lista as declarações Lean dos resultados desta seção, em \code{FormalizacaoAltura}, \code{FormalizacaoCrit}, \code{FormalizacaoCrit2} e \code{FormalizacaoCrit3}.

\begin{leantable}{Declarações Lean dos pisos, da necessidade e da altura crítica}{lp:tab:critical-lean}{Resultado}{Declarações Lean}
Quantização com salto limitado & \leanname{quantization_jump}, \leanname{nodal_tree_gen}, \leanname{nodal_tree_coarse} \\
Saltos nodais uniformes & \leanname{nodal_jump} \\
Escala do bump e piso geométrico & \leanname{besovNorm_bumpS_le}, \leanname{geometric_dist_lower}, \leanname{geometric_floor}, \leanname{L1_affine_lower}, \leanname{terminal_affine} \\
Piso por inclinação & \leanname{slope_floor}, \leanname{two_interval_pointwise}, \leanname{slope_dist_lower} \\
Necessidade conjunta & \leanname{geometric_necessity}, \leanname{slope_necessity}, \leanname{arithmetic_necessity}, \leanname{geometric_necessity_realP}, \leanname{slope_necessity_realP}, \leanname{arithmetic_necessity_realP} \\
Transferência & \leanname{bestError_le_realP}, \leanname{realP_le_bestError_add} \\
Altura linear crítica & \leanname{critical_approx}, \leanname{critical_level}, \leanname{critical_lq}, \leanname{critical_rate}, \leanname{critical_char_radius}, \leanname{critical_radius_ball_and_inclusion} \\
\end{leantable}

\subsection{Reprodução da verificação}

Na raiz do projeto Lean, a compilação e a auditoria são executadas por
\begin{verbatim}
lake build Formalizacao
lake env lean Formalizacao/Auditoria.lean
\end{verbatim}
A auditoria executa 201 consultas sobre 196 declarações distintas e registra somente \code{propext}, \code{Classical.choice} e \code{Quot.sound}. A saída está em \code{lean4/formalizacao/auditoria_axiomas.txt}. O mapa \code{mapas/MAPA_TEX_LEAN.md} aponta os nomes qualificados e os arquivos de origem. Na raiz do artigo, \texttt{python3 mapas/verificar\_mapa.py} confere os rótulos TeX e os nomes Lean.

\ifbesovabnt\else\articlereferences\fi
\end{document}
